\documentclass[11pt]{amsart}
\usepackage{amsmath,amssymb,amsthm,mathtools,mathrsfs}
\usepackage[margin=1.08in]{geometry}
\usepackage{booktabs}
\usepackage{array}
\usepackage{hyperref}
\hypersetup{hidelinks}
\usepackage[final,expansion=false]{microtype}
\setlength{\parskip}{4pt plus 1pt}
\emergencystretch=2.5em

\newtheorem{theorem}{Theorem}[section]
\newtheorem{proposition}[theorem]{Proposition}
\newtheorem{lemma}[theorem]{Lemma}
\newtheorem{corollary}[theorem]{Corollary}
\theoremstyle{definition}
\newtheorem{definition}[theorem]{Definition}
\newtheorem{problem}[theorem]{Open problem}
\theoremstyle{remark}
\newtheorem{remark}[theorem]{Remark}
\newtheorem{example}[theorem]{Example}

\newcommand{\C}{\mathbb C}
\newcommand{\R}{\mathbb R}
\newcommand{\Schatten}{\mathcal S}
\newcommand{\Ran}{\operatorname{Ran}}
\newcommand{\Spec}{\operatorname{Spec}}
\newcommand{\Tr}{\operatorname{Tr}}
\newcommand{\rank}{\operatorname{rank}}
\newcommand{\dist}{\operatorname{dist}}
\newcommand{\id}{\mathrm I}
\newcommand{\XiR}{\Xi}
\newcommand{\eps}{\varepsilon}
\newcommand{\dd}{\,d}

\title[Cofinal Weil quotients]
{Cofinal Weil Quotients:\\
Exact Transport, Fixed-Cutoff Summability,\\
and the Global Branch Obstruction}

\author{David Alejandro Trejo Pizzo}
\thanks{Independent Researcher. \texttt{dtrejopizzo@gmail.com}}

\begin{document}

\begin{abstract}
Connes--Consani--Moscovici construct finite self-adjoint operators from
principal Fourier restrictions of the semilocal Weil form, after shifting by
the least Ritz value and quotienting by the resulting radical.  Suzuki
develops a related continuous localized operator and formulates an expected
strong-resolvent limit.  Starting from these two operator programs, we
develop the analytic chain from finite radical quotients to a cutoff-free
spectral-gap formulation of Weil positivity.

At a fixed physical cutoff \(L\), we construct canonical isometries between
successive positive quotients.  The associated generator defect is computed
exactly.  Before the metric correction it is a sum of two rank-one channels
of opposite parity; the correction is the double-operator-integral transform
of a rank-two commutator.  We give both a divided-difference formula and an
explicit Loewner integral, then conjugate the calculation to the natural
Euclidean self-adjoint coordinates.  This yields exact resolvent and
new-mode-shell identities.

For the scalar second-resolvent trace we prove the characteristic formula
\[
 \Tr(s-D_N)^{-2}
 =\sum_{|k|\le N}(s-2\pi k/L)^{-2}
  -\bigl(\log\mathcal C_N(s)\bigr)'',
\]
where \(\mathcal C_N\) is the Cauchy transform of the least Ritz vector.  It
implies absolute local-uniform summability along a cofinal subsequence at
every fixed \(L\).  We also prove a uniform Weyl bound and show that
summability of the transported resolvent defect automatically implies
summability of the two-dimensional shell of new modes.  For consecutive
cutoffs we isolate one sufficient gap--drop series and exhibit a compact-form
counterexample showing that abstract Galerkin compactness alone cannot prove
that series.  We identify the fixed-\(L\) sum with the logarithmic curvature
of the selected continuum ground transform.  After an explicit exterior
mesh error \(O_K(L^2/N)\), the global limit is equivalent to

\[
 \partial_z^2\log(\widehat\phi_L/\widehat k_L)\longrightarrow0,
\]

where \(k_L\) is the CCM prolate model.

We then analyze this global gate.  Logarithmic-curvature convergence is
proved equivalent, modulo affine normalization, to transform convergence,
and a contour formula shows that it carries the complete zero-divisor
content.  Failure of RH is shown to produce fixed compactly supported
negative tests in both parity sectors.  Replacing the fixed model by a
moving co-Poisson radical yields an exact identity retaining primes, prime
powers, Gamma and pole jointly, and a constrained prolate ladder with
levels \(d_4\) and \(d_8\), \(d_4/d_8\asymp\lambda^{-8}\).

The branch choice is first isolated as a signed source estimate.  On
each principal-angle plane the trace projection has eigenvalues
\(\{-\chi,+\chi\}\): concentration determines its square but not its sign.
A translated Montgomery--Vaughan estimate reaches the required
\(O(\log\lambda)\) prime scale on each Fourier block and the correct local
\(O(\lambda^2)\) phase-space budget, but loses the block center and does not
select the global branch.  We then recover that center by an exact identity:
the actual prime-power--Gamma multiplier is a nonnegative full-line jump
form minus an explicit scalar \(\kappa_N=4\lambda+o(\lambda)\), and has its
unique minimum at \(t=0\).  In the even sector the pole is a positive
rank-one term.  A second exact identity cancels the continuous PNT main
measure against the growing pole branch and combines the remaining pole
branch with Gamma.  The sole unresolved arithmetic distribution is
\(d(\psi(e^u)-e^u)\), coupled to an explicit locally singular kernel.
The full diagonal becomes
\(QW_L=\mathcal E_*-c_*I-\mathcal A_\Delta\), with
\(\mathcal E_*\ge0\) and \(c_*>0\).  This gives a precise rank-one
spectral/discrepancy form of the branch gate and an unconditional
far-center theorem for Fejer packets.  An exact nonlocal Picone identity
and a positive one-atom countermodel show that generic jump positivity and
rank-one perturbation cannot prove the remaining one-sided estimate.

For comparison with the classical Weil criterion \cite{Weil1952}, we then pass from the
moving finite-cutoff branch problem to Riemann's full positive kernel
\(K\), normalized by \(\widehat K=\Xi\).  Smooth truncations
of \(K\) are operator quasimodes with super-exponentially small residual.
For every even compactly supported multiplier \(r\), the complete Weil form
has the exact cutoff-free representation
\[
 QW(Kr,Kr)=\mathscr E_K(r)-\frac12\operatorname{Var}_{\mu_K}(r),
\]
where \(\mathscr E_K\) contains every ordinary von Mangoldt atom and the
Gamma jump measure, and \(\mu_K\) is the probability measure induced by the
two polar evaluations.  Consequently the classical Weil equivalence takes
the sharp
Poincare inequality
\[
 \mathscr E_K(r)\ge\frac12\operatorname{Var}_{\mu_K}(r)
\]
on the even test core.  Each individual prime-power atom is load-bearing:
deleting even one of them makes this sharp inequality false on compact
truncations of an exact Riemann radical.  The exact theta dilation also shows
why retaining only divisible theta indices is insufficient: the
nondivisible indices and the central crossing region carry part of the sharp
constant.

This last equivalence is a normalized ground-state rewriting of Weil
positivity, not a proof or strengthening of it.  Analytic inequalities and
Wiener-space dynamics equivalent to RH were already developed by
Sun \cite{Sun}.  The finite-to-global construction developed on top of the CCM and
Suzuki frameworks therefore ends at two precisely separated tasks.  Full operator
convergence still requires the signed branch selector and the moving/model
identification.  A direct closure of RH requires the all-atom sharp
Doob--variance inequality above, equivalently the exclusion of spectrum
below \(1/2\) in the even mean-zero sector of the full Doob operator.  That
inequality is not proved here, and no conclusion about RH is claimed.
\end{abstract}

\maketitle
\tableofcontents

\section{Introduction and scope}

Put
\[
 \xi(s)=\tfrac12s(s-1)\pi^{-s/2}\Gamma(s/2)\zeta(s),
 \qquad \XiR(t)=\xi(\tfrac12+it).
\]
The spectral construction of Connes--Consani--Moscovici (CCM)
\cite{CCM} begins with the semilocal Weil form at a physical cutoff
\(L=2\log\lambda\), restricts it to a finite Fourier space, shifts by the
least Ritz value, and quotients by the one-dimensional radical.  The
resulting finite operator is self-adjoint in the positive quotient metric,
and its characteristic function is a finite real-rooted approximation.

There are two limiting operations:
\[
 N\longrightarrow\infty\quad(L\text{ fixed}),
 \qquad L\longrightarrow\infty.
\tag{1.1}
\]
The first asks whether the finite radical quotients and their arithmetic
generators form a coherent system.  The second asks whether the resulting
spectral divisor is the divisor of \(\XiR\).  The purpose of this paper is to
resolve a substantial part of the first question without conflating it with
the second.

The main conclusions are:
\begin{enumerate}
\item there is a canonical isometric connection between successive shifted
Weil quotients;
\item its arithmetic intertwining defect admits an exact rank-two/Loewner
formula;
\item scalar second-resolvent traces are absolutely summable along a cofinal
subsequence for every fixed \(L\);
\item the apparent new-mode shell is not an independent summability
obligation;
\item consecutive absolute summability requires source regularity not
provided by abstract compactness;
\item logarithmic curvature is an exact relative-divisor observable;
\item failure of RH is detected by fixed tests in both parity sectors;
\item a moving co-Poisson vector gives exact coupled
prime--Gamma--pole cancellation and a \(d_4/d_8\) prolate hierarchy;
\item the actual prime--Gamma multiplier has an exact centered jump-square
decomposition with unique Fourier center \(t=0\);
\item after exact PNT compensation, the missing global input is a one-sided
estimate for one explicit compensated Stieltjes functional, not three
separate source bounds or another absolute norm estimate.
\item a zero-extension Picone identity is exact but lossless, and a
positive one-atom countermodel rules out generic jump positivity as the
source of the branch sign.
\item smooth truncations of Riemann's full kernel are genuine operator
quasimodes, not only scalar near-radicals;
\item ordinary \(L^2\)-cyclicity cannot transfer that quasimode estimate to
the complementary floor at the required quantitative scale;
\item the cutoff-free completed Weil form is exactly a full-kernel
Doob energy minus one half of an explicit polar variance;
\item every literal von Mangoldt atom, every nondivisible theta remainder
and the central crossing region are load-bearing at that sharp constant.
\end{enumerate}

All finite-level statements below are conditional only on the same
simple-even and nonzero-source-overlap hypotheses needed to define the CCM
radical quotient in the indicated coordinate.  No location hypothesis on
the zeros of \(\zeta\) is used.

\section{Starting point in the CCM--Suzuki operator program}
\label{sec:boundary}

CCM construct the finite Weil matrices, identify their rank-two commutator
with the periodic derivative, shift by the least eigenvalue, quotient by the
radical, and construct the finite rank-one perturbation; see
\cite[Section~5, especially Lemma~5.4 and Theorem~5.10]{CCM}.  Their two
limiting parameters provide the finite spectral architecture used here:
first the Fourier dimension tends to infinity at fixed physical cutoff,
then the physical cutoff tends to infinity.

Suzuki constructs a continuous localized operator from the Weil
distribution and studies its finite-cutoff structure \cite{Suzuki}.  His
continuous screw-function coordinate and expected strong-resolvent limit
provide the second starting point.  Both programs seek a limiting
self-adjoint realization of the zero divisor; neither treats convergence as
an automatic consequence of the finite constructions.

The present paper builds the intermediate analytic layers needed to move in
that direction.  It first constructs a metric-compatible transport between
successive CCM quotients, computes its defect and proves fixed-cutoff
summability.  It then identifies the outer logarithmic-curvature gate,
tracks the two branch signs, and keeps primes, prime powers, Gamma and pole
coupled through the source compensation.  Finally it removes the cutoff
from the algebra by using Riemann's full positive kernel as an exact Weil
radical.  This last step converts the same completed Weil form into a sharp
Doob--variance problem.

The Doob coordinate is not asserted to have been the stated endpoint of
either source paper.  It is a direct ground-state coordinate for the Weil
form that underlies both operator programs.  It therefore separates two
questions which should not be conflated: convergence of the particular
finite operators to the zero divisor, and direct positivity of the limiting
completed Weil form.  The first retains the moving/model identification;
the second is equivalent to the sharp full-kernel inequality proved below
to be the sole remaining condition in that coordinate.

\section{Finite Weil data and the canonical connection}
\label{sec:finite}

Fix \(L>0\), put \(c_L=2\pi/L\), and let
\[
 V_N=\operatorname{span}\{U_k:|k|\le N\},
 \qquad D_{0,N}U_k=kU_k.
\tag{3.1}
\]
Let \(I_N:V_N\to V_{N+1}\) be Fourier inclusion.  The exact principal
restriction property of the finite Weil matrices is
\[
 W_N=I_N^*W_{N+1}I_N.
\tag{3.2}
\]
Their entries contain the coupled polar, Gamma and prime-power terms of the
semilocal Weil form.

Assume that the least eigenvalue is simple, its eigenvector is even, and its
overlap with
\(\eta_N=\sum_{|k|\le N}U_k\) is nonzero.  Normalize the ground vector by
\[
 \eps_N=\min\Spec W_N,
 \qquad T_N=W_N-\eps_N\id,
 \qquad \langle\eta_N,\xi_N\rangle=1.
\tag{3.3}
\]
Define
\[
 K_N=(\C\xi_N)^\perp,\quad P_N=P_{K_N},\quad
 S_N=T_N|_{K_N},\quad
 a_N=D_{0,N}\xi_N,\quad b_N=P_N\eta_N.
\tag{3.4}
\]
The physical quotient operator represented on \(K_N\) is
\[
 D_N=c_L\bigl(P_ND_{0,N}P_N-|a_N\rangle\langle b_N|\bigr),
 \qquad S_ND_N=D_N^*S_N.
\tag{3.5}
\]

Set
\[
 \Delta_N=\eps_N-\eps_{N+1}\ge0,
 \qquad C_N=P_{N+1}I_N|_{K_N}.
\tag{3.6}
\]

\begin{lemma}[Shifted compression]
For \(x,y\in K_N\),
\[
 \langle S_{N+1}C_Nx,C_Ny\rangle
 =\langle(S_N+\Delta_N\id)x,y\rangle.
\tag{3.7}
\]
\end{lemma}

\begin{proof}
Equation (3.2) gives, for \(x,y\in V_N\),
\[
 \langle T_{N+1}I_Nx,I_Ny\rangle
 =\langle T_Nx,y\rangle+\Delta_N\langle x,y\rangle.
\]
Projecting \(I_Nx\) and \(I_Ny\) away from \(\ker T_{N+1}\) does not alter
the left side, which proves (3.7).
\end{proof}

Put
\[
 B_N=(S_N+\Delta_N\id)^{-1/2}S_N^{1/2},
 \qquad j_N=C_NB_N.
\tag{3.8}
\]

\begin{theorem}[Canonical radical connection]
The map \(j_N:K_N\to K_{N+1}\) is injective and
\[
 j_N^*S_{N+1}j_N=S_N.
\tag{3.9}
\]
Thus it is an isometry for the positive shifted-Weil quotient metrics.
\end{theorem}

\begin{proof}
The right side of (3.7) is positive for \(x\ne0\), so \(C_N\) is
injective.  Substitution of (3.8) into (3.7), using functional calculus for
the commuting functions of \(S_N\), gives (3.9).
\end{proof}

The uncorrected inclusion cannot itself descend to the quotient unless
\(\Delta_N=0\): if \(T_N\xi_N=0\), then
\(\|I_N\xi_N\|_{T_{N+1}}^2=\Delta_N\|\xi_N\|^2\).  Formula (3.8) is exactly
the correction of this obstruction.

\section{Exact cross-level generator defect}
\label{sec:defect}

Define
\[
 e_N=I_N\xi_N-\xi_{N+1},\qquad
 u_N=D_{0,N+1}e_N,\qquad
 v_N=P_{N+1}I_N\xi_N=P_{N+1}e_N.
\tag{4.1}
\]

\begin{theorem}[Two-channel uncorrected defect]
On \(K_N\),
\[
\boxed{
 X_N:=D_{N+1}C_N-C_ND_N
 =c_L|u_N\rangle\langle b_N|
  +\frac{c_L}{\|\xi_N\|^2}|v_N\rangle\langle a_N|.}
\tag{4.2}
\]
The first channel maps even vectors to odd vectors and the second maps odd
vectors to even vectors.  Hence their nonzero singular values separate:
\[
 c_L\|u_N\|\,\|b_N\|,
 \qquad
 \frac{c_L\|v_N\|\,\|a_N\|}{\|\xi_N\|^2}.
\tag{4.3}
\]
Moreover,
\[
\boxed{\|S_{N+1}^{1/2}v_N\|^2
       =\Delta_N\|\xi_N\|^2.}
\tag{4.4}
\]
\end{theorem}

\begin{proof}
Write \(A'_N=D_{0,N}-|a_N\rangle\langle\eta_N|\), so
\(A'_N\xi_N=0\).  For \(x\in K_N\), projection away from the new ground
line in the argument of \(A'_{N+1}\) may be omitted.  Add and subtract
\(I_NA'_Nx\):
\[
\begin{aligned}
X_Nx/c_L
={}&P_{N+1}(A'_{N+1}I_N-I_NA'_N)x\\
 &+P_{N+1}I_N(\id-P_N)A'_Nx.
\end{aligned}
\]
The first line is \(|u_N\rangle\langle b_N|x\).  Since evenness makes
\(a_N\) odd and \((A'_N)^*\xi_N=a_N\), the second is
\(\|\xi_N\|^{-2}|v_N\rangle\langle a_N|x\).  Orthogonality of the parity
sectors proves (4.3).

For (4.4), apply the shifted compression identity before restricting to
\(K_N\) to \(\xi_N\), and note that replacing \(I_N\xi_N\) by \(v_N\)
changes it by a vector in \(\ker T_{N+1}\).
\end{proof}

\section{The metric correction as a Loewner transform}
\label{sec:loewner}

From (3.5),
\[
 D_N-D_N^*
 =c_L\bigl(|b_N\rangle\langle a_N|-|a_N\rangle\langle b_N|\bigr),
\tag{5.1}
\]
and metric self-adjointness gives
\[
\boxed{
[D_N,S_N]
=c_L\bigl(|b_N\rangle\langle S_Na_N|
          -|a_N\rangle\langle S_Nb_N|\bigr).}
\tag{5.2}
\]
This commutator has rank at most two.

Let
\[
 f_\Delta(t)=\sqrt{\frac{t}{t+\Delta}},
 \qquad B_N=f_{\Delta_N}(S_N).
\tag{5.3}
\]
If \(S_N=\sum_p\sigma_pE_p\), the finite-dimensional double-operator-
integral identity is
\[
\boxed{
[D_N,B_N]
=\sum_{p,q}f_{\Delta_N}^{[1]}(\sigma_p,\sigma_q)
 E_p[D_N,S_N]E_q,}
\tag{5.4}
\]
where \(f^{[1]}(x,y)=(f(y)-f(x))/(y-x)\), with the derivative on the
diagonal.  The representation
\[
 f_\Delta(t)=\frac1\pi\int_0^1
 \frac{t}{t+r\Delta}\frac{\dd r}{\sqrt{r(1-r)}}
\tag{5.5}
\]
and (5.2) give the more explicit formula
\[
\boxed{
\begin{aligned}
[D_N,B_N]
=\frac{c_L\Delta_N}{\pi}\int_0^1\sqrt{\frac r{1-r}}
\bigl(&|R_rb_N\rangle\langle S_NR_ra_N|\\
      &-|R_ra_N\rangle\langle S_NR_rb_N|\bigr)\dd r,
\end{aligned}}
\tag{5.6}
\]
with \(R_r=(S_N+r\Delta_N\id)^{-1}\).

Since \(f_\Delta\) is operator monotone, its Loewner matrix is positive.
If \(m_N=\min\Spec S_N\), the induced completely positive Schur multiplier
gives, for every \(1\le p\le\infty\),
\[
\boxed{
 \|[D_N,B_N]\|_{\Schatten_p}
 \le
 \frac{\Delta_N}
 {2\sqrt{m_N}(m_N+\Delta_N)^{3/2}}
 \|[D_N,S_N]\|_{\Schatten_p}.}
\tag{5.7}
\]

\begin{theorem}[Complete transport defect]
The arithmetic intertwining defect
\(\mathcal R_N=D_{N+1}j_N-j_ND_N\) is
\[
\boxed{\mathcal R_N=X_NB_N+C_N[D_N,B_N].}
\tag{5.8}
\]
Consequently it is determined explicitly by (4.2) and (5.6), with no input
from the zeros of \(\XiR\).
\end{theorem}

\begin{proof}
Use \(j_N=C_NB_N\), add and subtract \(C_ND_NB_N\), and apply (4.2).
\end{proof}

Raw Euclidean norms of \(\mathcal R_N\) do not measure the canonical
transport.  Put
\[
 \mathsf A_N=S_N^{1/2}D_NS_N^{-1/2},\qquad
 \mathsf U_N=S_{N+1}^{1/2}C_N(S_N+\Delta_N\id)^{-1/2}.
\tag{5.9}
\]
Then \(\mathsf A_N\) is Euclidean self-adjoint and \(\mathsf U_N\) is an
isometry.  The correct defect is
\[
 \widehat{\mathcal R}_N
 =\mathsf A_{N+1}\mathsf U_N-\mathsf U_N\mathsf A_N
 =S_{N+1}^{1/2}\mathcal R_NS_N^{-1/2}.
\tag{5.10}
\]
Writing \(H_N=S_N+\Delta_N\id\),
\[
\boxed{
 \widehat{\mathcal R}_N
 =S_{N+1}^{1/2}X_NH_N^{-1/2}+\mathsf U_NY_N,}
\tag{5.11}
\]
where
\[
\begin{aligned}
Y_N=\frac{c_L\Delta_N}{\pi}\int_0^1\sqrt{\frac r{1-r}}
\bigl(&|H_N^{1/2}R_rb_N\rangle\langle S_N^{1/2}R_ra_N|\\
      -&|H_N^{1/2}R_ra_N\rangle\langle S_N^{1/2}R_rb_N|\bigr)\dd r.
\end{aligned}
\tag{5.12}
\]
The two singular values of the first term in (5.11) are
\[
 \alpha_N=c_L\|S_{N+1}^{1/2}u_N\|\,\|H_N^{-1/2}b_N\|,
\quad
 \beta_N=\frac{c_L\sqrt{\Delta_N}}{\|\xi_N\|}
          \|H_N^{-1/2}a_N\|.
\tag{5.13}
\]

\section{Resolvent transport and the two-mode shell}
\label{sec:resolvent}

For \(z\notin\R\), let
\[
 G_N(z)=(\mathsf A_N-z)^{-1},\qquad
 E_N=G_{N+1}\mathsf U_N-\mathsf U_NG_N.
\tag{6.1}
\]
The resolvent identity gives
\[
\boxed{E_N=-G_{N+1}\widehat{\mathcal R}_NG_N.}
\tag{6.2}
\]
Let
\[
 Q_N^{\rm new}=\id-\mathsf U_N\mathsf U_N^*.
\tag{6.3}
\]
The dimension of its range is two, one mode in each parity sector.

\begin{proposition}[Exact second-resolvent increment]
\[
\boxed{
\begin{aligned}
\Tr G_{N+1}(z)^2-\Tr G_N(z)^2
={}&\Tr\mathsf U_N^*(G_{N+1}E_N+E_NG_N)\\
 &+\Tr Q_N^{\rm new}G_{N+1}(z)^2Q_N^{\rm new}.
\end{aligned}}
\tag{6.4}
\]
\end{proposition}

\begin{proof}
Insert \(\id=\mathsf U_N\mathsf U_N^*+Q_N^{\rm new}\) into the trace of
\(G_{N+1}^2\).  In the transported block use
\(G_{N+1}\mathsf U_N=\mathsf U_NG_N+E_N\), expand once from each side, and
use cyclicity of the finite trace.
\end{proof}

There is also an exact Hilbert--Schmidt identity.  Put
\(H_N(z)=\|G_N(z)\|_{\Schatten_2}^2\).  Then
\[
\boxed{
\begin{aligned}
\|G_{N+1}Q_N^{\rm new}\|_{\Schatten_2}^2
={}&H_{N+1}-H_N
-2\Re\Tr((\mathsf U_NG_N)^*E_N)-\|E_N\|_{\Schatten_2}^2.
\end{aligned}}
\tag{6.5}
\]

\section{The scalar characteristic and fixed-\texorpdfstring{\(L\)}{L}
summability}
\label{sec:summability}

Let \(\phi_N\) be an \(L^2\)-unit least Ritz vector of \(W_N\).  Define
\[
 \mathcal C_N(s)=\sum_{k=-N}^{N}\frac{(\phi_N)_k}{s-c_Lk},
 \qquad s\in\C\setminus\R,
\tag{7.1}
\]
and
\[
 \tau_N(s)=\Tr(s-D_N)^{-2}.
\tag{7.2}
\]

\begin{theorem}[Exact Cauchy-transform formula]
\[
\boxed{
 \tau_N(s)=\sum_{k=-N}^{N}\frac{1}{(s-c_Lk)^2}
            -\bigl(\log\mathcal C_N(s)\bigr)''.}
\tag{7.3}
\]
\end{theorem}

\begin{proof}
First take the dimensionless diagonal \(D_0e_k=ke_k\), normalize
\(\xi_N\) by \(\langle\eta_N,\xi_N\rangle=1\), and write
\(D'_N=D_0-|D_0\xi_N\rangle\langle\eta_N|\).  The matrix determinant
lemma gives
\[
\begin{aligned}
\det(s-D'_N)
&=\prod_{k=-N}^{N}(s-k)
  \left(1+\sum_{k=-N}^{N}\frac{k(\xi_N)_k}{s-k}\right)\\
&=s\prod_{k=-N}^{N}(s-k)
  \sum_{k=-N}^{N}\frac{(\xi_N)_k}{s-k}.
\end{aligned}
\tag{7.4}
\]
The factor \(s\) is the radical eigenvalue and is removed on the quotient.
Rescale \(s\mapsto s/c_L\); multiplication of the ground vector by a
nonzero scalar disappears after two logarithmic derivatives.  Finally use
\(\Tr(s-D)^{-2}=-(\log\det(s-D))''\).
\end{proof}

Finite self-adjointness implies that \(\mathcal C_N\) is zero-free off the
real axis.

\subsection{Quotient trace versus the full regularized determinant}

Formula (7.3) is the trace of the finite quotient.  The full CCM
regularized determinant also retains the free exterior Fourier mesh.  In the
present affine coordinate its second-logarithmic-derivative contribution is
\[
 T^{\rm ext}_{L,N}(s)
 =\sum_{|k|>N}\frac1{(s-c_Lk)^2}.
\tag{7.5}
\]
Thus, up to the sign convention used for logarithmic curvature and an
affine regularization phase (which disappears after two derivatives), the
full trace is \(\tau_{L,N}+T^{\rm ext}_{L,N}\).  If
\(K\Subset\C\setminus\R\) and \(N\) is large enough relative to \(L\) and
\(K\), then
\[
 \sup_{s\in K}|T^{\rm ext}_{L,N}(s)|
 \le C_K\frac{L^2}{N}.
\tag{7.6}
\]
Indeed, \(|s-c_Lk|\gg_K |k|/L\) on the exterior tail and
\(\sum_{k>N}k^{-2}=O(N^{-1})\).  Therefore fixed-\(L\) passage to
\(N\to\infty\) is harmless.  On a diagonal limit, however, discarding the
exterior factor requires
\[
 \frac{N(L)}{L^2}\longrightarrow\infty.
\tag{7.7}
\]
No such rate is required if the known exterior factor is retained or
divided out exactly.  This normalization distinction is essential: the
finite quotient trace and the full regularized determinant are not the same
object before the exterior mesh is accounted for.

\subsection{A sufficient consecutive rate}

Let
\[
 g_{N+1}=\lambda_2(W_{N+1})-\eps_{N+1}>0,
\tag{7.8}
\]
and choose the phases so that
\(\langle I_N\phi_N,\phi_{N+1}\rangle\ge0\).

\begin{lemma}[Ground-line rotation]
\[
 \|I_N\phi_N-\phi_{N+1}\|^2
 \le \frac{2\Delta_N}{g_{N+1}}.
\tag{7.9}
\]
\end{lemma}

\begin{proof}
Write \(I_N\phi_N=a\phi_{N+1}+h\), with \(a\ge0\) and
\(h\perp\phi_{N+1}\).  The min--max principle yields
\(\Delta_N\ge g_{N+1}\|h\|^2\).  Since \(a^2+\|h\|^2=1\) and
\(1-a\le1-a^2\), one has
\(\|I_N\phi_N-\phi_{N+1}\|^2=2(1-a)\le2\|h\|^2\).
\end{proof}

\begin{theorem}[Consecutive scalar criterion]
If
\[
 \sum_{N\ge N_0}\sqrt{\Delta_N/g_{N+1}}<\infty,
\tag{7.10}
\]
then, for every compact \(K\Subset\C\setminus\R\),
\[
 \sum_{N\ge N_0}\sup_{s\in K}
 |\tau_{N+1}(s)-\tau_N(s)|<\infty.
\tag{7.11}
\]
\end{theorem}

\begin{proof}
For \(j=0,1,2\), the vectors
\[
 \left(\frac{d^j}{ds^j}\frac1{s-c_Lk}\right)_{k\in\mathbb Z}
\]
belong uniformly to \(\ell^2(\mathbb Z)\) for \(s\in K\).  Equations
(7.9)--(7.10) imply
\(\sum_N\|\mathcal C_{N+1}-\mathcal C_N\|_{C^2(K)}<\infty\).
The limit is nonzero and, by Hurwitz, zero-free off \(\R\).  After finitely
many indices, \(f\mapsto(\log f)''\) is Lipschitz on this bounded subset of
\(C^2(K)\).  The two new lattice summands in (7.3) are \(O_K(N^{-2})\).
Summing proves (7.11).
\end{proof}

\subsection{Rate-free cofinal summability}

\begin{theorem}[Fixed-\(L\) cofinal summability]
Assume that the quotient construction is defined on a cofinal set of
levels.  There exists a cofinal sequence \(N_j\to\infty\) such that, for
every compact \(K\Subset\C\setminus\R\),
\[
\boxed{
 \sum_j\sup_{s\in K}
 |\tau_{N_{j+1}}(s)-\tau_{N_j}(s)|<\infty.}
\tag{7.12}
\]
The same sequence can be chosen simultaneously for a countable exhaustion
of \(\C\setminus\R\).
\end{theorem}

\begin{proof}
The unit Ritz vectors have uniformly bounded shifted form norm because
their Rayleigh values decrease to the lower bound of the closed semilocal
Weil form.  The compact form embedding proved in the CCM construction makes
them precompact in \(L^2\).  Select coherent phases and a subsequence with
\[
 \phi_{N_j}\to\phi,\qquad
 \|\phi_{N_j}-\phi\|_2\le2^{-j-2}.
\tag{7.13}
\]
Lower semicontinuity shows that \(\phi\) is a continuum ground state and
has norm one.  Hence
\(\sum_j\|\phi_{N_{j+1}}-\phi_{N_j}\|_2<\infty\).

The \(\ell^2\) estimate used above now gives bounded variation of
\(\mathcal C_{N_j}\) in \(C^2(K)\).  Its limit cannot vanish identically,
because the Cauchy transform determines all Fourier coefficients of the
nonzero vector \(\phi\).  Hurwitz gives zero-freeness in both half-planes.
Finally, the lattice shells
\[
 \sum_{N_j<|k|\le N_{j+1}}(s-c_Lk)^{-2}
\]
are disjoint and absolutely summable, uniformly on \(K\).  Formula (7.3)
proves (7.12).  A diagonal selection over a countable compact exhaustion
gives the last assertion.
\end{proof}

\begin{corollary}[Selected continuum transform]
The Fourier transform of the continuum ground state selected in the proof
has no nonreal zero.
\end{corollary}

\begin{proof}
Compact support and \(L^2\)-convergence imply locally uniform convergence of
the Fourier transforms.  Each finite transform has only real zeros by the
CCM determinant theorem.  The limit is not identically zero, so Hurwitz
excludes nonreal zeros.
\end{proof}

\begin{theorem}[Identification of the fixed-\(L\) sum]
Let \(\phi_L\) be the continuum ground state selected above and write its
Fourier transform in the CCM normalization as
\[
 \widehat\phi_{L,N}(z)
 =2L^{-1/2}\sin\left(\frac{Lz}{2}\right)
   \sum_{k=-N}^{N}\frac{(\phi_{L,N})_k}{z-c_Lk}.
\tag{7.14}
\]
Then the limit in (7.12) is not merely an abstract holomorphic function:
\[
\boxed{
 \tau_L(z):=\lim_{j\to\infty}\tau_{L,N_j}(z)
 =-\partial_z^2\log\widehat\phi_L(z)}
\tag{7.15}
\]
locally uniformly on \(\C\setminus\R\).
\end{theorem}

\begin{proof}
Compact support turns \(L^2\)-convergence of the selected ground vectors
into locally uniform convergence of their Fourier transforms with all
derivatives.  The full regularized determinant is, up to a nonzero scalar,
\[
 F_{L,N}(z)=-ie^{-iLz/2}\widehat\phi_{L,N}(z).
\]
The exponential is affine after taking a logarithm, hence invisible to the
second derivative.  The full and quotient traces differ by the exterior
tail (7.5), which tends locally uniformly to zero at fixed \(L\) by (7.6).
Taking the limit in
\((\log F_{L,N})''=-\tau^{\rm full}_{L,N}\) proves (7.15).
\end{proof}

\section{Uniform Weyl control and elimination of the shell}
\label{sec:shell}

\begin{theorem}[Uniform fixed-\(L\) Weyl bound]
For each \(z\notin\R\),
\[
\boxed{\sup_N\|G_N(z)\|_{\Schatten_2}^2<\infty.}
\tag{8.1}
\]
\end{theorem}

\begin{proof}
Let \(\mathcal K_L\) be the \(L^2\)-closure of the unit Ritz vectors.  The
compact form embedding makes \(\mathcal K_L\) compact and disjoint from
zero.  Fourier transform is continuous and injective on this set.  Hence
there exist finitely many points \(w_1,\ldots,w_r\) and \(c>0\) such that
\[
 \max_j|\widehat f(w_j)|\ge c
 \qquad(f\in\mathcal K_L).
\tag{8.2}
\]
After translating the support to \([-L/2,L/2]\),
\[
 |\widehat f(w)|\le\sqrt L\,e^{(L/2)|\Im w|}.
\tag{8.3}
\]
Jensen's formula centered at a point satisfying (8.2), together with
(8.3), gives the uniform zero count
\[
 \#\{\theta\in\Spec\mathsf A_N:|\theta|\le R\}
 \le C_L(1+R).
\tag{8.4}
\]
Here multiplicity is retained through the finite characteristic identity.
Dyadic summation of (8.4) yields
\[
 \sup_N\sum_{\theta\in\Spec\mathsf A_N}|\theta-z|^{-2}<\infty,
\]
which is (8.1).
\end{proof}

\begin{corollary}[The new shell is not independent]
If
\[
 \sum_N\|E_N\|_{\Schatten_1}
 =\sum_N\|G_{N+1}\widehat{\mathcal R}_NG_N\|_{\Schatten_1}<\infty,
\tag{8.5}
\]
then
\[
\boxed{
 \sum_N\|G_{N+1}Q_N^{\rm new}\|_{\Schatten_2}^2<\infty.}
\tag{8.6}
\]
\end{corollary}

\begin{proof}
Since \(\|\mathsf U_NG_N\|\le|\Im z|^{-1}\),
\[
 |\Tr((\mathsf U_NG_N)^*E_N)|
 \le|\Im z|^{-1}\|E_N\|_{\Schatten_1}.
\]
Hypothesis (8.5) also implies
\(\sum_N\|E_N\|_{\Schatten_2}^2<\infty\).  Sum the nonnegative left side
of (6.5) and use the uniform bound (8.1).
\end{proof}

Thus the two-series trace-Cauchy package reduces to the single transported
series (8.5).

\section{Why compactness does not prove consecutive summability}
\label{sec:counterexample}

CCM's monotone Ritz theorem gives
\[
 \sum_N\Delta_N
 =\eps_{N_0}-\inf\Spec A_L<\infty,
\tag{9.1}
\]
but this does not imply (7.10), even when the limiting spectral gap is
positive.

\begin{proposition}[Compact-form counterexample]
There is a closed positive form with compact form embedding, a simple ground
state, a positive limiting gap and nested coordinate cores for which
\[
 \Delta_N\asymp N^{-2},\qquad g_N\ge g>0,\qquad
 \sum_N\sqrt{\Delta_N/g_N}=\infty.
\tag{9.2}
\]
\end{proposition}

\begin{proof}
On \(\mathcal H=\C\oplus\ell^2(\mathbb N)\), define
\[
 q(x)=\sum_{n\ge1}n\,|x_n-n^{-3/2}x_0|^2.
\tag{9.3}
\]
The \(x_0\)-direction is finite-dimensional, and on \(x_0=0\) the form
controls \(\sum n|x_n|^2\); hence its form embedding is compact.  Its kernel
is the simple line generated by \((1,n^{-3/2})\), and compactness gives a
positive gap above it.

Restrict to coordinates \(0,1,\ldots,N\).  Eliminating the last \(N\)
coordinates in the least-eigenvalue equations gives exactly
\[
 t_N=\eps_N\left(1+\sum_{n\le N}\frac{1}{n^2(n-\eps_N)}\right),
 \qquad t_N=\sum_{n>N}\frac{1}{n^2}.
\tag{9.4}
\]
Therefore
\[
 \eps_N\sim\frac{1}{(1+\zeta(3))N},
 \qquad
 \Delta_N\sim\frac{1}{(1+\zeta(3))N^2},
\tag{9.5}
\]
which proves (9.2).
\end{proof}

Consequently, consecutive summability must use additional regularity of the
actual polar--Gamma--Euler ground state.  For example, a source-specific
Ritz estimate
\[
 \eps_N-\inf\Spec A_L=O(N^{-p}\log N),\qquad p>1,
\tag{9.6}
\]
together with a limiting positive gap would suffice.  Neither audited source
proves such a rate.

\section{Reproducible finite-level check}

The accompanying NumPy-only script
\texttt{106\_04\_cofinal\_weil\_defect\_probe.py}, retained in the research
archive, constructs one largest Weil matrix and obtains all lower levels by
principal compression.  It checks the shifted compression, the rank-two
formula (4.2), the complete defect (5.8), the metric conjugation (5.10), the
resolvent identity (6.2), and the shell increment (6.4).  A separate scalar
script checks (7.3) and the summability bookkeeping.  These calculations are
diagnostic: floating-point decay is not used in any theorem.

\section{The remaining global limit}
\label{sec:global}

Put
\[
 s=\frac12+iz,
 \qquad \XiR(z)=\xi\!\left(\frac12+iz\right).
\tag{11.1}
\]
The chain rule gives the sign-correct identity
\[
\boxed{
 \partial_z^2\log\XiR(z)
 =-\left(\frac{\xi'}{\xi}\right)'\!\left(\frac12+iz\right).}
\tag{11.2}
\]
If \(\Theta_{L,N}^{\rm full}=\frac12+iD_{L,N}^{\rm full}\), then
\[
 \Tr(s-\Theta_{L,N}^{\rm full})^{-2}
 =-\tau_{L,N}^{\rm full}(z).
\tag{11.3}
\]
Thus the desired global statement can be written equivalently as
\[
\begin{aligned}
 \Tr(s-\Theta_{L,N(L)}^{\rm full})^{-2}
 &\longrightarrow-\left(\frac{\xi'}{\xi}\right)'(s),\\
 \tau_{L,N(L)}^{\rm full}(z)
 &\longrightarrow
 \left(\frac{\xi'}{\xi}\right)'\!\left(\frac12+iz\right).
\end{aligned}
\tag{11.4}
\]
There is no additional sign or scale factor.

Let \(k_L\) be the explicit prolate model denoted \(k_\lambda\) by CCM,
where \(L=2\log\lambda\).  Their model theorem gives
\[
 \widehat k_L\longrightarrow\XiR
\tag{11.5}
\]
locally uniformly on closed substrips of \(|\Im z|<1/2\).  On compact sets
which avoid the zeros of \(\XiR\), Cauchy's theorem therefore gives
\[
 \partial_z^2\log\widehat k_L
 \longrightarrow\partial_z^2\log\XiR.
\tag{11.6}
\]
Combining (7.15) and (11.6) yields an exact affine-free reduction.

\begin{theorem}[Affine-free ground/model criterion]
Assume that a simple-even continuum ground state \(\phi_L\) is selected for
cofinally many \(L\).  Choose the finite diagonal \(N=N(L)\) from the
fixed-\(L\) cofinal sequences so that, on a compact exhaustion,
\[
 \tau_{L,N(L)}^{\rm q}-\tau_L\longrightarrow0
 \quad\text{locally uniformly}.
\tag{11.7}
\]
Such a diagonal exists by (7.12).  If the exterior factor is discarded,
choose it simultaneously so that \(N(L)/L^2\to\infty\); if that factor is
retained or divided out exactly, no exterior rate is needed.  Under these
diagonal conditions, the global second-resolvent identification (11.4)
holds if and only if
\[
\boxed{
 \partial_z^2\log
 \left(\frac{\widehat\phi_L(z)}{\widehat k_L(z)}\right)
 \longrightarrow0}
\tag{11.8}
\]
locally uniformly on every simply connected compact set in the open strip
which avoids the zeros of both transforms.
\end{theorem}

\begin{proof}
Choose compatible logarithms on the compact set.  The identity
\[
\begin{aligned}
 \partial_z^2\log\widehat\phi_L-\partial_z^2\log\XiR
={}&\partial_z^2\log\frac{\widehat\phi_L}{\widehat k_L}
 +\partial_z^2\log\frac{\widehat k_L}{\XiR}
\end{aligned}
\tag{11.9}
\]
is exact.  Its last term tends to zero by (11.5) and Cauchy estimates.
Now use (7.15) and (11.2)--(11.4).  The finite-diagonal assertion follows
from the exterior bound (7.6).
\end{proof}

The curvature criterion is invariant under
\(\widehat\phi_L(z)\mapsto e^{a_L+b_Lz}\widehat\phi_L(z)\); hence it removes
the full admissible affine normalization.

\begin{lemma}[Quantitative relative-error transfer]
Let \(K\) be compact and let \(K_r\) be its closed \(r\)-neighborhood
inside a zero-free domain.  Put
\[
 m_{L,K,r}=\inf_{z\in K_r}|\widehat k_L(z)|,
 \qquad
 \eta_{L,K,r}=\sup_{z\in K_r}
 |c_L^*\widehat\phi_L(z)-\widehat k_L(z)|.
\tag{11.10}
\]
If \(\eta_{L,K,r}\le m_{L,K,r}/2\), then
\[
\boxed{
 \sup_{z\in K}\left|
 \partial_z^2\log\frac{c_L^*\widehat\phi_L(z)}{\widehat k_L(z)}
 \right|
 \le\frac4{r^2}\frac{\eta_{L,K,r}}{m_{L,K,r}}.}
\tag{11.11}
\]
Here \(c_L^*\ne0\) is any source normalization.
\end{lemma}

\begin{proof}
Write \(c_L^*\widehat\phi_L/\widehat k_L=1+h_L\).  On \(K_r\),
\(|h_L|\le1/2\) and \(|\log(1+h_L)|\le2|h_L|\).  Cauchy's estimate for
the second derivative on disks of radius \(r\) gives (11.11).
\end{proof}

A stronger source-side sufficient condition is
\[
 \sup_{|\beta|\le B}
 \int_{\lambda^{-1}}^\lambda
 |c_L^*\phi_L(u)-k_L(u)|u^\beta\,d^*u
 \longrightarrow0,
 \qquad B<\frac12.
\tag{11.12}
\]
Indeed, \(|u^{-iz}|=u^{\Im z}\), so (11.12) bounds the numerator in
(11.10) on closed substrips.  Model convergence keeps the denominator
uniformly positive away from the zeros of \(\XiR\), and the lemma implies
(11.8).

In the operator program adopted here, (11.8) and (11.12) are the
force-bearing global source-identification estimates.  They do not follow
from fixed-\(L\) compactness, and the present paper does not assume them.

\section{Curvature rigidity and the zero divisor}
\label{sec:rigidity}

The remaining limit is not weakened by taking two logarithmic derivatives.
It is merely made insensitive to affine normalization.

For a zero-free holomorphic function \(F\), put
\[
 \mathcal K_F:=-(\log F)''.
\]

\begin{theorem}[Affine rigidity]
\label{thm:affine-rigidity}
Let \(\Omega\subset\C\) be simply connected, and let \(F_j,G_j\) be
holomorphic and zero-free on \(\Omega\).  Then
\[
 \mathcal K_{F_j}-\mathcal K_{G_j}\longrightarrow0
\]
locally uniformly if and only if there are \(a_j,b_j\in\C\) such that
\[
 e^{-a_j-b_jz}\frac{F_j(z)}{G_j(z)}\longrightarrow1
\]
locally uniformly.  If \(0\in\Omega\), both functions are even, and
\(F_j(0)=G_j(0)\ne0\), then \(a_j=b_j=0\).
\end{theorem}

\begin{proof}
Choose \(H_j=\log(F_j/G_j)\), fix \(z_0\in\Omega\), and subtract
\[
 A_j(z)=H_j(z_0)+(z-z_0)H_j'(z_0).
\]
Two integrations along bounded paths in a slightly larger compact set give
\[
 \sup_K|H_j-A_j|\le C_K\sup_{K'}|H_j''|.
\]
This proves one direction after exponentiation; the converse follows from
Cauchy estimates.  In the even normalized case, choose \(H_j(0)=0\).
Then \(H_j'(0)=0\), so the affine correction vanishes.
\end{proof}

Curvature also counts the relative zero divisor directly.  If \(F\) is
holomorphic near the closure of a Jordan domain \(D\), with no zero on
\(C=\partial D\), then for every \(z_*\in\C\),
\[
 \boxed{
 N_F(D)=-\frac1{2\pi i}\int_C(z-z_*)\mathcal K_F(z)\,dz.}
 \tag{12.1}
\]
Indeed, a zero of multiplicity \(m\) contributes
\(\mathcal K_F(z)=-m(z-a)^{-2}+O(1)\), and the residue in (12.1) is
\(-m\).

\begin{theorem}[Divisor rigidity]
\label{thm:divisor-rigidity}
Assume that every selected transform \(\widehat\phi_L\) has only real
zeros, \(\widehat k_L\to\XiR\) locally uniformly, and
\[
 \mathcal K_{\widehat\phi_L}-\mathcal K_{\widehat k_L}\longrightarrow0
\]
uniformly on every fixed contour avoiding both divisors.  Then every zero
of \(\XiR\) is real.
\end{theorem}

\begin{proof}
If \(\XiR\) had a nonreal zero, choose a small disk \(D\) around it,
disjoint from the real axis.  Rouch\'e's theorem gives
\(N_{\widehat k_L}(D)>0\) for large \(L\), whereas real-rootedness gives
\(N_{\widehat\phi_L}(D)=0\).  Formula (12.1) and curvature convergence
force their difference to tend to zero, a contradiction.
\end{proof}

Thus (11.8) already has the full zero-location content of RH.  The next
sections isolate its exact variational and arithmetic requirements.

\section{Variational selection and bilateral detection}
\label{sec:selection}

Let \(A_L\) be the even continuum semilocal Weil operator, with normalized
simple ground state \(\phi_L\), first two eigenvalues
\(\epsilon_{0,L}<\epsilon_{1,L}\), and gap
\(g_L=\epsilon_{1,L}-\epsilon_{0,L}\).  For a normalized model \(q_L\),
put \(R_L=\langle A_Lq_L,q_L\rangle\).  Spectral expansion gives
\[
 1-|\langle q_L,\phi_L\rangle|^2
 \le\frac{R_L-\epsilon_{0,L}}{g_L}.
 \tag{13.1}
\]
Consequently, with \(L=2\log\lambda\),
\[
 \sup_{|\Im z|\le B}
 |\widehat q_L(z)-e^{i\vartheta_L}\widehat\phi_L(z)|
 \le W_{L,B}
 \left(\frac{2(R_L-\epsilon_{0,L})}{g_L}\right)^{1/2},
 \tag{13.2}
\]
where
\[
 W_{L,B}^2=\frac{\lambda^{2B}-\lambda^{-2B}}{2B},
 \qquad W_{L,0}^2=L.
\]
Together with (11.11), this proves the curvature limit whenever
\[
 \boxed{
 \frac{\|k_L\|_2W_{L,B}}{\inf|\widehat k_L|}
 \left(\frac{R_L-\epsilon_{0,L}}{g_L}\right)^{1/2}
 \longrightarrow0.}
 \tag{13.3}
\]

There is a weaker gate which proves RH without identifying the complete
ground transform.

\begin{theorem}[Residual/overlap criterion]
\label{thm:residual-overlap}
Let \(A_L\phi_L=\epsilon_L\phi_L\), \(\|\phi_L\|=\|q_L\|=1\), and
\(\mu_L\in\R\).  Then
\[
 \boxed{
 |\epsilon_L-\mu_L|\,|\langle\phi_L,q_L\rangle|
 \le\|(A_L-\mu_L)q_L\|.}
 \tag{13.4}
\]
Hence \(\mu_L\to0\) and a vanishing residual/overlap quotient force
\(\epsilon_L\to0\).
\end{theorem}

\begin{proof}
Take the inner product of \((A_L-\mu_L)q_L\) with \(\phi_L\), use
self-adjointness, and apply Cauchy--Schwarz.
\end{proof}

To use (13.4), failure of RH must be detected in the selected parity.  The
following strengthens the usual unspecified-parity Weil test.

\begin{theorem}[Bilateral parity detection]
\label{thm:bilateral}
If RH is false, there exist fixed nonzero real functions
\[
 f_+\in C_c^\infty(\R)^{\rm even},
 \qquad f_-\in C_c^\infty(\R)^{\rm odd},
\]
such that
\[
 QW(f_+,f_+)<0,
 \qquad QW(f_-,f_-)<0.
 \tag{13.5}
\]
Consequently both semilocal parity bottoms are bounded above by a fixed
negative constant for every sufficiently large support window.
\end{theorem}

\begin{proof}
Use the centered zero parameter \(s\), so that an off-line quartet is
\(\{s_0,\bar s_0,-s_0,-\bar s_0\}\), with
\(\Re s_0\Im s_0\ne0\).  For a real even test with transform \(F\), the
quartet contributes \(4m\Re F(s_0)^2\); for a real odd test it contributes
\(-4m\Re F(s_0)^2\).  Thus the target values \(F_+(s_0)=i\) and
\(F_-(s_0)=1\) make both contributions negative.

Choose a real even Paley--Wiener damping transform \(\Phi\), first fixing
it and then a finite zero box on whose complement
\(|\Phi(s)/\Phi(s_0)|\le1/2\).  Delete the full target orbit, with
multiplicity, and let \(Z_R\) be the real even polynomial annihilating the
remaining finite divisor.  Since \(\Im(s_0^2)\ne0\), the real-linear map
\((a,b)\mapsto a+bs_0^2\) is an isomorphism.  Hence one can choose real
coefficients so that
\[
 F_{N,+}(z)=Z_R(z)(a_N+b_Nz^2)\Phi(z)^N,
\]
\[
 F_{N,-}(z)=izZ_R(z)(c_N+d_Nz^2)\Phi(z)^N
\]
have the prescribed target values.  Paley--Wiener gives fixed-parity
compactly supported smooth inverse transforms.  Rapid strip decay and
the zero count \(N(T)=O(T\log T)\) dominate the remaining zero sum by a
summable majorant times \(4^{-N}\).  For one sufficiently large finite
\(N\), both full Weil values are negative.  Rayleigh's principle proves
the semilocal assertion.
\end{proof}

Combining Theorems~\ref{thm:residual-overlap} and
\ref{thm:bilateral}, either parity ground proves RH if
\[
 \boxed{
 \frac{\|A_Lq_L\|}{|\langle\phi_L,q_L\rangle|}\longrightarrow0.}
 \tag{13.6}
\]
This gate is weaker than (13.3), but neither an endpoint normalization nor
qualitative positivity can control its denominator.

\begin{remark}[Sharpness of both quotients]
For \(0<\eta<1\), the positivity-improving matrix
\[
 A_\eta=\begin{pmatrix}-1&-\eta\\-\eta&0\end{pmatrix},
 \qquad q_\eta=\frac{(\eta^2,1)}{\|(\eta^2,1)\|}
\]
has a simple strictly positive ground whose eigenvalue tends to \(-1\),
while \(\|A_\eta q_\eta\|\sim\eta\) and its overlap with the ground is also
asymptotic to \(\eta\).  Thus the quotient in (13.6) does not improve.

Likewise, on \(L^2([-1,1])^{\rm even}\), the two real-rooted transforms
\(\sin z/z\) and \((\sin(z/2)/(z/2))^2\) can be made respectively the ground
and a quasimode with Rayleigh excess tending to zero but a gap collapsing at
the same rate.  Their logarithmic curvatures remain different.  Hence
compact resolvent, simplicity, real-rootedness and an unscaled quasimode
estimate cannot replace (13.3).
\end{remark}

\section{Near-radical clusters and the model scale}
\label{sec:cluster}

Let \(K\) be Riemann's full additive kernel, normalized by
\(\widehat K=\XiR\).  Every derivative \(K_j=\partial_x^jK\) lies in the
polarized Weil radical:
\[
 QW(K_j,g)=0,
 \tag{14.1}
\]
because \(\widehat K_j(z)=(iz)^j\XiR(z)\) vanishes on the complete
nontrivial divisor.

For \(b>1/2\), define
\[
 \|f\|_{\mathcal X_b}
 =\sum_{r=0}^2\sup_x
 e^{b|x|}(1+|x|)^3|f^{(r)}(x)|.
\]
The completed source formula satisfies the unconditional continuity bound
\[
 \boxed{|QW(f,g)|\le C_b\|f\|_{\mathcal X_b}\|g\|_{\mathcal X_b}.}
 \tag{14.2}
\]
Indeed, two derivatives control the symmetric Gamma difference at the
origin, exponential decay controls its tail and the polar evaluations, and
the prime powers are bounded by
\[
 \sum_{m\ge2}\frac{\Lambda(m)}{m^{1/2+b'}}
 =-\frac{\zeta'}\zeta\left(\frac12+b'\right)<\infty
 \qquad(1/2<b'<b).
\]

Choose a smooth cutoff \(\chi_a\), equal to one on \([-a,a]\), and put
\(v_{j,a}=\chi_aK_j\), \(e_{j,a}=v_{j,a}-K_j\), with
\(a=\log\lambda\).  The theta-series decay of \(K\) gives
\[
 \|e_{j,a}\|_{\mathcal X_b}
 \le C_j\lambda^{M_j}e^{-c_j\lambda^2}.
\]
Using (14.1) twice,
\[
 QW(v_{i,a},v_{j,a})=QW(e_{i,a},e_{j,a}).
\]
Therefore, for \(i,j\in\{0,1,2\}\),
\[
 \boxed{
 |QW(v_{i,a},v_{j,a})|
 \le C\lambda^Me^{-c\lambda^2}.}
 \tag{14.3}
\]

Min--max applied to \(\operatorname{span}\{v_{0,a},v_{2,a}\}\) and
\(\operatorname{span}\{v_{1,a}\}\) gives unconditional upper bounds of
this size for the second even and first odd eigenvalues.  Under RH, and
only then using Weil positivity, the next-even and parity gaps collapse at
the same double-exponential-in-\(L\) scale.  Thus no fixed, polynomial or
ordinary exponential gap can close (13.3).

The quantitative CCM model estimate has a different and much coarser
scale.  With \(k_\lambda\) the compact prolate model and \(k\) the full
kernel, the Meixner estimate and the Gaussian exterior tail give
\[
 \boxed{
 \|k_\lambda-k\|_2^2=O(\lambda^{-1}),
 \qquad
 \sup_{|\Im z|\le B}|\widehat k_\lambda(z)-\XiR(z)|
 =O_B(\lambda^{B-1/2})}
 \tag{14.4}
\]
for every fixed \(B<1/2\).

This does not control the Weil form.  The Gamma multiplier grows like
\((2\pi)^{-1}\log|t|\), and hence is unbounded in the \(L^2\) topology.
The finite prime block satisfies only
\[
 \left\|\sum_{1<n\le\lambda^2}\Lambda(n)T(n)\right\|
 \le4C_\psi\lambda,
 \tag{14.5}
\]
which consumes the full squared gain in (14.4).  Although radicality gives
the exact scalar identity
\[
 QW(k_\lambda,k_\lambda)=QW(k_\lambda-k,k_\lambda-k),
 \tag{14.6}
\]
the available model estimate proves neither \(R_L\to0\), nor an operator
residual, nor the weighted quotient (13.3).

\section{The moving co-Poisson radical}
\label{sec:moving}

The fixed model error in (14.4) can be removed from the exact identity by
moving inside the full co-Poisson radical.

For an even compactly supported BV function \(f\), supported in
\([-\lambda,\lambda]\) and satisfying \(\int_0^\infty f=0\), define
\[
 \mathcal E f(u)=u^{1/2}\sum_{n\ge1}f(nu).
 \tag{15.1}
\]
Then \(F=\mathcal E f\in L^2(\R_+^*,d^*u)\), and throughout
\(\Re s>-1/2\),
\[
 \boxed{
 \mathcal MF(s)=\zeta(s+\tfrac12)\mathcal Mf(s+\tfrac12).}
 \tag{15.2}
\]

\begin{proof}[Proof of (15.2)]
Partition the half-line into \(((n-1)u,nu]\).  Zero mass and bounded
variation give
\[
 \left|u\sum_{n\ge1}f(nu)\right|
 \le u\operatorname {Var}(f),
\]
which proves the \(L^2\) bound at zero; compact support gives the upper
cutoff.  Absolute summation proves (15.2) for \(\Re s>1/2\).  The pole at
\(s=1/2\) is removed by zero mass, and analytic continuation gives the
asserted half-plane.
\end{proof}

In particular, the Mellin transform of \(F\) vanishes at both coordinates
attached to every nontrivial zero.  Let \(P_I\) be restriction to
\(I_\lambda=[\lambda^{-1},\lambda]\), and put
\[
 V_\lambda=P_IF,
 \qquad R_\lambda=(1-P_I)F.
\]

\begin{theorem}[Moving-radical cancellation]
\label{thm:moving-radical}
For every semilocal core vector \(G\) supported in \(I_\lambda\),
\[
 \boxed{
 QW_\lambda(V_\lambda,G)=-QW(R_\lambda,G),
 \qquad
 QW_\lambda(V_\lambda,V_\lambda)=QW(R_\lambda,R_\lambda).}
 \tag{15.3}
\]
The right sides are absolutely convergent zero-side values.
\end{theorem}

\begin{proof}
At every zero, the transform of \(R_\lambda\) is the negative of that of
\(V_\lambda\).  Substitute this equality into the polarized zero-side
formula, first with \(G\), and then in both entries.
\end{proof}

If \(H=V_\lambda^**G\), Connes' semilocal trace formula
\cite{Connes2026} rewrites (15.3)
without separating any source term:
\[
\boxed{
 W_{0,2}(H)+\log(TW)H(1)
 +\Tr\!\left(\vartheta(H)
 (1-P_T^S-\widehat P_W^S)\right)
 =-QW(R_\lambda,G).}
\tag{15.4}
\]
This is the exact joint prime--prime-power--Gamma--pole cancellation used
below.

Poisson summation gives, for \(0<u<\lambda^{-1}\),
\[
 R_\lambda(u)
 =\mathcal E((1-P_\lambda)\widehat f)(u^{-1})
 -\tfrac12u^{1/2}f(0).
 \tag{15.5}
\]
The standard midpoint convention is understood for zero-extended BV
functions.  Imposing also \(f(0)=0\) leaves only Fourier leakage.

Let \(\psi_{j,\lambda}\), \(j=0,4,8,12\), be normalized even prolate
functions in the Fourier \(+1\) sector, with compressed eigenvalues
\(\chi_j\), and put \(d_j=1-\chi_j^2\).  The fixed-order asymptotics and
compressed-Fourier identities are those used in \cite{CCM}.  The cofactor
combination of the
first three modes annihilates both rows
\[
 (\psi_j(0))_j,
 \qquad (\chi_j\psi_j(0))_j,
\]
and therefore satisfies exactly \(f(0)=\int f=0\).  On the two-dimensional
four-mode constrained space, the angle form
\[
 \mathfrak A_\lambda(f,g)
 =\langle(1-P_\lambda)\widehat f,
          (1-P_\lambda)\widehat g\rangle
\]
has generalized eigenvalues
\[
 \boxed{
 \alpha_{0,\lambda}\asymp d_4,
 \qquad
 \alpha_{1,\lambda}\asymp d_8,
 \qquad
 \frac{d_4}{d_8}\asymp\lambda^{-8},}
 \tag{15.6}
\]
where
\[
 d_4\asymp\lambda^9e^{-4\pi\lambda^2},
 \qquad d_8\asymp\lambda^{17}e^{-4\pi\lambda^2}.
 \tag{15.7}
\]

To verify (15.6), put \(q_r=1-\chi_{4r}\).  The constraints in the
coordinates \(x_r=\psi_{4r}(0)c_r\) are
\(\sum x_r=\sum q_rx_r=0\), while
\(q_r/q_{r+1}=O(\lambda^{-8})\).  The min--max test vectors
\[
 (q_1-q_2,q_2-q_0,q_0-q_1,0),
 \quad
 (0,q_2-q_3,q_3-q_1,q_1-q_2)
\]
give the matching upper bounds \(q_1,q_2\); the two constraints give the
corresponding lower bounds.  Since \(d_j\asymp q_j\), (15.6) follows.

The polynomial separation in (15.6) is more than sufficient for the
Paley--Wiener weight in (13.3).  The unresolved issue is transferring this
angle hierarchy to the signed Weil operator and selecting the correct
branch.

\section{Why abstract geometry does not select that branch}
\label{sec:nogos}

Three natural selectors fail for exact reasons.

First, the completed semigroup is not positivity preserving.  In the
additive coordinate, for disjoint real tests \(f,g\), the off-diagonal pole
and Gamma kernels are respectively
\[
 2\cosh\frac{x-y}{2},
 \qquad
 -c_\infty(|x-y|),
 \quad
 c_\infty(u)=\frac{e^{-u/2}}{1-e^{-2u}}.
\]
The prime block is supported only when \(|x-y|=\log p^k\).  Choose two
nonnegative disjoint bumps whose separation is near \(1\), avoiding every
prime-power logarithm.  Since
\[
 2\cosh(u/2)-c_\infty(u)>0
\]
there, their cross form is positive.  This violates the first
Beurling--Deny criterion \(QW_L(f,g)\le0\) for positivity preservation.
The construction can be symmetrized, so it also rules out the natural
pointwise cone in the even sector.

Second, adjoining finitely many derivatives of the full kernel does not
repair the branch.  Every transform in
\[
 \operatorname{span}\{K,K',\ldots,K^{(m)}\}
\]
contains the common factor \(\XiR\).  A hypothetical off-line evaluation
mode is therefore asymptotically orthogonal to every such fixed family.
The same observation applies to every finite subspace of the complete
co-Poisson radical.

Nor does real-rootedness prevent escape.  If \(g_\delta\) is the normalized
indicator of \([-\delta,\delta]\), then
\[
 \psi_a(x)=2^{-1/2}\bigl(g_\delta(x-a)+g_\delta(x+a)\bigr)
\]
is positive, even and normalized, and
\(\widehat\psi_a(z)=\sqrt2\cos(az)\widehat g_\delta(z)\) has only real
zeros.  Nevertheless \(\psi_a\rightharpoonup0\) and becomes orthogonal to
every fixed radical model as \(a\to\infty\).

The exact many-vector obstruction is as follows.

\begin{theorem}[Frame--residual conservation]
\label{thm:frame-residual}
Let \(A\) be lower-semibounded, self-adjoint and of compact resolvent, and
let \(P_-={\bf1}_{(-\infty,0)}(A)\).  For
\(q_1,\ldots,q_M\in\operatorname{Dom}(A)\), define
\(Qc=\sum c_jq_j\), \(D=AQ\), and
\[
 a_-(Q)=\inf_{\substack{v\in\Ran P_-\\\|v\|=1}}\|Q^*v\|^2.
\]
If \(a_-(Q)>0\), then
\[
 \Spec(A|_{\Ran P_-})
 \subset\left[-\frac{\|D^*P_-\|}{\sqrt{a_-(Q)}},0\right),
 \tag{16.1}
\]
and for an orthonormal negative eigenbasis \((\phi_\ell)\),
\[
 \boxed{
 \sum_{j=1}^M\|P_-Aq_j\|^2
 =\sum_{\epsilon_\ell<0}\epsilon_\ell^2
   \sum_{j=1}^M|\langle\phi_\ell,q_j\rangle|^2
 \ge a_-(Q)\sum_{\epsilon_\ell<0}\epsilon_\ell^2.}
 \tag{16.2}
\]
\end{theorem}

\begin{proof}
The negative space reduces \(A\).  Apply the lower frame inequality to
\(Av\) and use \(Q^*Av=D^*v\) to obtain (16.1).  Expanding each
\(P_-Aq_j\) in the negative eigenbasis and applying Parseval proves
(16.2).
\end{proof}

Thus adding modes multiplies the useful overlap and the projected residual
by the same amount.  A fixed negative branch forces their quotient not to
vanish.

There is an exact full-radical version.  If \(S\) is any finite-dimensional
subspace of the polarized Weil radical and \(QW(h,h)<0\), then
\[
 h_S=h-P_Sh\perp S,
 \qquad QW(h_S,h_S)=QW(h,h)<0.
\]
Thus no finite radical frame removes a negative direction; truncation can
see it only through the aggregate exterior defect.

There is also a dimension mismatch in the presently available coercivity.
Write \(A_L=G_L+B_L\), where
\[
 \langle G_Lf,f\rangle
 \ge c_0\int\log(2+|t|)|\widehat f(t)|^2\frac{dt}{2\pi}
      -C_0\|f\|^2
\]
and \(\|B_L\|\le b_L\).  If
\(T_L=\exp(2(C_0+b_L)/c_0)\), a negative unit vector has at least half of
its Fourier mass in \([-T_L,T_L]\).  The time--frequency trace bound then
gives
\[
 \dim\Ran P_-\le\frac{2LT_L}{\pi}.
 \tag{16.3}
\]
The absolute prime--pole estimate \(b_L=O(\lambda)\) yields only
\(O(Le^{C\lambda})\) possible negative modes, whereas the Slepian family
has dimension \(O(\lambda^2)\).  A signed estimate at logarithmic effective
scale is therefore necessary.

\section{The signed prime--Gamma branch gate}
\label{sec:spg}

The trace projection itself contains an unavoidable sign choice.  Let
\(P,Q\) be orthogonal projections, and let a principal vector satisfy
\(PQP e=\chi^2e\).  With \(s=(1-\chi^2)^{1/2}\), the principal plane has
\[
 P=\begin{pmatrix}1&0\\0&0\end{pmatrix},
 \qquad
 Q=\begin{pmatrix}\chi^2&\chi s\\\chi s&s^2\end{pmatrix}.
\]
Therefore, for \(\mathcal J=I-P-Q\),
\[
 \boxed{
 \mathcal J^2=\chi^2I,
 \qquad \Spec\mathcal J=\{-\chi,+\chi\}.}
 \tag{17.1}
\]
When the leakage \(d=1-\chi^2\) tends to zero, both branches remain of
order one.  Prolate concentration fixes the square in (17.1), not the sign
of its square root.

Let \(q_L^{\rm mov}=V_\lambda/\|V_\lambda\|\) be the normalized first
moving vector and \(e_L^{\rm ext}=R_\lambda/\|V_\lambda\|\).  Put
\[
 r_L^{\rm mov}=\langle A_Lq_L^{\rm mov},q_L^{\rm mov}\rangle,
 \quad
 r_L^\perp=(1-|q_L^{\rm mov}\rangle\langle q_L^{\rm mov}|)
             A_Lq_L^{\rm mov}.
\]
The moving-radical identity gives the exact first-variation formulas
\[
 \boxed{
 \|r_L^\perp\|
 =\sup_{\substack{g\perp q_L^{\rm mov}\\\|g\|=1}}
 |QW(e_L^{\rm ext},g)|,
 \qquad
 r_L^{\rm mov}=QW(e_L^{\rm ext},e_L^{\rm ext}).}
 \tag{17.2}
\]
The residual is linear in the exterior error, while the Rayleigh value is
quadratic.

Even if one adds a polynomial logarithmic graph estimate, termwise control
gives at best
\[
 \|r_L^\perp\|
 =O(\operatorname{poly}(\lambda)\sqrt{d_4}).
\]
The branch scale required below is \(o(\lambda^{-\sigma}d_8)\),
\(\sigma<1/2\), and their ratio contains
\[
 \frac{\sqrt{d_4}}{\lambda^{-\sigma}d_8}
 \asymp\lambda^{\sigma-25/2}e^{2\pi\lambda^2}.
 \tag{17.3}
\]
Thus no polynomial regularity estimate reaches the selector: the complete
linear variation must cancel.

Define the complementary bottom
\[
 \gamma_L^\perp
 =\inf_{\substack{g\perp q_L^{\rm mov}\\\|g\|=1}}
   QW_L(g,g).
\]
If \(r_L^{\rm mov}<\gamma_L^\perp\), scalar Schur complementation gives
\[
 0\le r_L^{\rm mov}-\epsilon_{0,L}
 \le\frac{\|r_L^\perp\|^2}
          {\gamma_L^\perp-r_L^{\rm mov}},
 \qquad
 \epsilon_{1,L}\ge\gamma_L^\perp.
 \tag{17.4}
\]

\begin{problem}[Gate SPG]
\label{prob:spg}
Prove, from the actual ordinary-prime source, that
\[
 r_L^{\rm mov}=O(d_4),
 \qquad
 \gamma_L^\perp\ge c d_8,
 \qquad
 \lambda^\sigma
 \frac{\|r_L^\perp\|}
      {\gamma_L^\perp-r_L^{\rm mov}}
 \longrightarrow0
 \quad(\sigma<1/2).
 \tag{17.5}
\]
Equivalently, for \(\mathcal H=(q_L^{\rm mov})^**g\), keep all terms
coupled and prove
\[
 \sup_{\substack{g\perp q_L^{\rm mov}\\\|g\|=1}}
 \left|
 W_{0,2}(\mathcal H)+\log(TW)\mathcal H(1)
 +\Tr\!\left(\vartheta(\mathcal H)
   (1-P_T^S-\widehat P_W^S)\right)
 \right|
 =o(\lambda^{-\sigma}d_8),
 \tag{17.6}
\]
and the corresponding diagonal lower bound \(cd_8\) on the orthogonal
complement.
\end{problem}

Gate SPG and Theorem~\ref{thm:bilateral} imply RH through
Theorem~\ref{thm:residual-overlap}, without identifying the entire ground
transform.  To deduce the stronger curvature limit (11.8), one additionally
needs the distinct moving/model identification
\[
 \partial_z^2\log
 \frac{\widehat q_L^{\rm mov}}{\widehat k_L}\longrightarrow0.
 \tag{MI}
\]
Neither implication is built into the other.

\section{A local logarithmic theorem and the lost center}
\label{sec:local-log}

The logarithmic effective scale in Section~\ref{sec:nogos} can be reached
locally.  Put
\[
 X_L=\lambda^2=e^L,
 \qquad
 \mathcal P_{X_L}(t)
 =\sum_{n\le X_L}\frac{\Lambda(n)}{\sqrt n}\,n^{-it}.
\]

\begin{theorem}[Translated prime mean square]
\label{thm:mv}
Uniformly for \(U\in\R\) and \(Y\ge1\),
\[
 \boxed{
 \int_{U-Y}^{U+Y}|\mathcal P_{X_L}(t)|^2dt
 \le C(YL^2+X_LL).}
 \tag{18.1}
\]
In particular, for \(Y=X_L/L\), the root-mean-square size is \(O(L)\).
\end{theorem}

\begin{proof}
The translated Montgomery--Vaughan generalized Hilbert inequality
\cite{MV} gives, for \(a_n=\Lambda(n)n^{-1/2-iU}\),
\[
 \int_{-Y}^{Y}\left|\sum_{n\le X_L}a_nn^{-it}\right|^2dt
 \le C\left(Y\sum|a_n|^2+\sum n|a_n|^2\right).
\]
Chebyshev bounds and partial summation give
\[
 \sum_{n\le X_L}\frac{\Lambda(n)^2}{n}=O(L^2),
 \qquad
 \sum_{n\le X_L}\Lambda(n)^2=O(X_LL),
\]
which proves (18.1).
\end{proof}

If \(v_1,\ldots,v_d\) are orthonormal and supported in the additive
interval of length \(L\), Bessel's inequality gives
\(\sum_j|\widehat v_j(t)|^2\le L\).  On a block of length \(X_L/L\),
Cauchy--Schwarz and (18.1) therefore yield
\[
 \sum_j\int_B(2\Re\mathcal P_{X_L})_+
       |\widehat v_j|^2\frac{dt}{2\pi}
 \le CLX_L.
 \tag{18.2}
\]
Near \(|t|\asymp X_L/L\), the Gamma multiplier is at least \(cL\).
Hence an orthonormal family with a fixed fraction of its mass in that block
and negative local prime--Gamma form has size \(O(X_L)=O(\lambda^2)\).
This is exactly the local Slepian budget.

The estimate does not globalize.  It is uniform in \(U\), and the absolute
prime norm confines possible negative wells only to
\(|t|\le e^{C\lambda}\).  Summing the local budget over all block centers
returns \(O(Le^{C\lambda})\).  Nor can (18.1) exclude one negative well.

The loss is certified by twisting the coefficients:
\[
 a_n^{(\tau)}=\frac{\Lambda(n)}{\sqrt n}n^{i\tau},
 \qquad
 \mathcal P_{X_L}^{(\tau)}(t)=\mathcal P_{X_L}(t-\tau).
\]
All mean-square estimates remain unchanged, while the coherent peak
\(\sum_{n\le X_L}\Lambda(n)/\sqrt n\asymp2\lambda\) moves to the arbitrary
center \(\tau\).  Thus magnitude estimates cannot select the fundamental
branch.  Gate SPG must use the untwisted locations and signs of the actual
von Mangoldt phases jointly with Gamma and the polar block.

\section{The centered prime--Gamma square}
\label{sec:centered-square}

The center discarded by Theorem~\ref{thm:mv} can be restored exactly.  Set
\[
 I_L=[-L/2,L/2],
 \qquad N=e^L=\lambda^2,
 \qquad w_n=\frac{\Lambda(n)}{\sqrt n},
 \qquad a_n=\log n.
\]
Extend \(F\in L^2(I_L)\) by zero to \(\R\), and let
\((\tau_aF)(x)=F(x-a)\).  This is full-line translation after zero
extension, not periodic translation on \(I_L\).

With
\[
 \theta(t)=
 \operatorname{Im}\log\Gamma\!\left(\frac14+\frac{it}{2}\right)
 -\frac t2\log\pi,
\]
the CCM normalization gives
\[
\begin{aligned}
 Q_L(F,F)
 &=
 \int_\R|\widehat F(t)|^2
 \left(
 2\theta'(t)-2\sum_{n\le N}w_n\cos(ta_n)
 \right)\frac{dt}{2\pi}\\
 &\quad+
 2\operatorname{Re}\!\left(
 \widehat F(i/2)\overline{\widehat F(-i/2)}
 \right).
\end{aligned}
\tag{19.1}
\]

Define the positive form
\[
\begin{aligned}
 \mathcal D_N(F,G)
 &:=
 \sum_{n\le N}w_n
 \langle F-\tau_{a_n}F,G-\tau_{a_n}G\rangle\\
 &\quad+
 \sum_{k=0}^{\infty}\int_0^\infty
 e^{-2(k+1/4)x}
 \langle F-\tau_xF,G-\tau_xG\rangle\,dx
\end{aligned}
\tag{19.2}
\]
and the explicit scalar
\[
 \kappa_N
 =
 2\sum_{n\le N}\frac{\Lambda(n)}{\sqrt n}
 -2\theta'(0).
\tag{19.3}
\]

\begin{theorem}[Exact centered square]
\label{thm:centered-square}
On the CCM logarithmic Gamma-form domain,
\[
\boxed{
\begin{aligned}
 Q_L(F,F)
 &=
 \mathcal D_N(F,F)-\kappa_N\|F\|_2^2\\
 &\quad+
 2\operatorname{Re}\!\left(
 \widehat F(i/2)\overline{\widehat F(-i/2)}
 \right).
\end{aligned}}
\tag{19.4}
\]
Moreover,
\[
\boxed{
 \kappa_N
 =
 2\sum_{n\le N}\frac{\Lambda(n)}{\sqrt n}
 +\gamma+\frac\pi2+3\log2+\log\pi
 =4\lambda+o(\lambda).}
\tag{19.5}
\]
\end{theorem}

\begin{proof}
Full-line Plancherel gives
\[
 \|F-\tau_aF\|_2^2
 =
 2\int_\R(1-\cos(at))|\widehat F(t)|^2\frac{dt}{2\pi}.
\tag{19.6}
\]
Writing \(b_k=k+1/4\), the digamma expansion gives
\[
\begin{aligned}
 2\bigl(\theta'(t)-\theta'(0)\bigr)
 &=
 \sum_{k=0}^{\infty}
 \frac{(t/2)^2}{b_k\bigl(b_k^2+(t/2)^2\bigr)}\\
 &=
 \sum_{k=0}^{\infty}
 2\int_0^\infty e^{-2b_kx}(1-\cos(tx))\,dx .
\end{aligned}
\tag{19.7}
\]
Hence the Fourier multiplier of \(\mathcal D_N\) is
\[
 2\sum_{n\le N}w_n(1-\cos(ta_n))
 +2(\theta'(t)-\theta'(0)).
\]
Subtracting (19.3) gives the multiplier in (19.1), which proves
(19.4).  Finally,
\[
 2\theta'(0)=\psi(1/4)-\log\pi
 =-\gamma-\frac\pi2-3\log2-\log\pi,
\]
and PNT with partial summation gives
\(\sum_{n\le N}\Lambda(n)/\sqrt n=2\sqrt N+o(\sqrt N)\).
\end{proof}

\begin{corollary}[Unique fundamental center]
\label{cor:unique-center}
For
\[
 m_N(t)=2\theta'(t)-2\sum_{n\le N}w_n\cos(t\log n),
\]
one has
\[
\boxed{
\begin{aligned}
 m_N(t)-m_N(0)
 &=
 2\sum_{n\le N}w_n(1-\cos(t\log n))\\
 &\quad+2(\theta'(t)-\theta'(0))\ge0,
\end{aligned}}
\tag{19.8}
\]
with strict inequality for \(t\ne0\).
\end{corollary}

\begin{proof}
Every summand is nonnegative, and the Gamma series in (19.7) is strictly
positive away from zero.
\end{proof}

Thus the artificial coefficient twist in Section~\ref{sec:local-log} is
excluded by the exact source identity: it moves the prime well without
moving the Gamma minimum.

There is also a pure-gauge representation.  Put
\[
 \phi_N(t)=
 \theta(t)-\sum_{n\le N}
 \frac{\Lambda(n)}{\sqrt n\log n}\sin(t\log n),
 \quad U_N=e^{i\phi_N},
 \quad\mathsf X=-i\partial_t.
\]
Then, on the smooth core,
\[
 2(U_N^*\mathsf XU_N-\mathsf X)=M_{m_N}.
\tag{19.9}
\]
This identity does not remove the finite-interval boundary or the polar
channel.

In the even sector,
\[
 \widehat F(i/2)=\widehat F(-i/2)
 =\langle h_L,F\rangle,
 \qquad h_L=\mathbf1_{I_L}\cosh(x/2),
\]
and therefore
\[
\boxed{
 Q_L^+(F,G)
 =
 \mathcal D_N(F,G)-\kappa_N\langle F,G\rangle
 +2\overline{\langle h_L,F\rangle}\langle h_L,G\rangle.}
\tag{19.10}
\]
If \(\mathcal L_N^+\) is associated with
\[
 \ell_N^+(F,G)
 =
 \mathcal D_N(F,G)
 +2\overline{\langle h_L,F\rangle}\langle h_L,G\rangle,
\]
then
\[
\boxed{A_L^+=\mathcal L_N^+-\kappa_NI,\qquad\mathcal L_N^+\ge0.}
\tag{19.11}
\]

Consequently Gate SPG can be restated as a quantitative rank-one spectral
theorem for the explicit positive operator \(\mathcal L_N^+\).  For the moving vector
\(q_L=q_L^{\rm mov}\), its cross part is
\[
\sup_{\substack{g\perp q_L\\\|g\|=1}}
\left|
\mathcal D_N(q_L,g)
+2\overline{\langle h_L,q_L\rangle}\langle h_L,g\rangle
\right|
=o(\lambda^{-\sigma}d_8),
\quad \sigma<\frac12,
\tag{19.12}
\]
and its complementary part is
\[
\inf_{\substack{g\perp q_L\\\|g\|=1}}
\left(\mathcal D_N(g,g)+2|\langle h_L,g\rangle|^2\right)
\ge\kappa_N+cd_8.
\tag{19.13}
\]
If the separate Gate-SPG Rayleigh estimate
\(r_L^{\rm mov}=O(d_4)\) is proved, then
\(\ell_N^+(q_L,q_L)=\kappa_N+O(d_4)\).  Together with that estimate,
(19.12)--(19.13) assert that \(q_L\) is a first-level
\(\mathcal L_N^+\)-quasimode to the required form-dual accuracy and that
the next level clears the explicit threshold.  Positivity of
\(\mathcal L_N^+\) alone does not imply these assertions.

For qualitative nonnegativity there is a narrower conditional rank-one
criterion.  Assume the force-bearing inertia statement
\[
 \lambda_1(\mathcal D_N)<\kappa_N<\lambda_2(\mathcal D_N).
\]
Then rank-one inertia gives
\[
 Q_L^+\ge0
 \quad\Longleftrightarrow\quad
 1+2\langle h_L,(\mathcal D_N-\kappa_N)^{-1}h_L\rangle\le0.
\tag{19.14}
\]
At finite Galerkin level the second factor is the determinant ratio
\[
 \frac{\det(\mathcal L_N^+-z)}
 {\det(\mathcal D_N-z)}
 =
 1+2\langle h_L,(\mathcal D_N-z)^{-1}h_L\rangle.
\tag{19.15}
\]
The sign in (19.14), rather than existence of the square in (19.4), is
arithmetic spectral information.  Criterion (19.14) concerns only
qualitative nonnegativity under the additional inertia hypothesis; it does
not imply the quantitative residual and \(d_8\) gap in
(19.12)--(19.13).

\begin{remark}[Form domain and boundary]
The Gamma jump density in (19.2) behaves as \(dx/(2x)\) near zero.  Define
the logarithmic form domain as the completion of \(C_c^\infty(I_L)\) under
\[
 \|F\|_2^2+
 \int_{\mathbb R}\log(2+|t|)|\widehat F(t)|^2\,dt.
\]
Then the identity extends to that domain by closure.
For translations comparable with \(L\), zero extension creates a boundary
loss which cancels the \(O(\lambda)\) constant in (19.3).  Replacing
\(\tau_a\) by periodic translation gives a false formula.
\end{remark}

\section{Exact PNT-discrepancy compensation}
\label{sec:pnt-compensation}

The same cancellation has a complementary Stieltjes form.  Let \(f,g\) be
smooth, zero-extended functions on \(I_L\), and put
\[
 H(u)=\langle f,\tau_ug\rangle,
 \qquad F(u)=H(u)+H(-u),
 \qquad \Psi(u)=\psi(e^u).
\tag{20.1}
\]
In this section \(Q_L=QW_L\) denotes the same semilocal CCM form.
For a Gate-SPG cross term \(g\perp f\), one has \(F(0)=0\).

The polar, Euler and off-diagonal Gamma terms are respectively
\[
\begin{aligned}
 P(f,g)
 &=\int_0^L F(u)(e^{u/2}+e^{-u/2})\,du,\\
 E(f,g)
 &=-\int_{(0,L]}F(u)e^{-u/2}\,d\Psi(u),\\
 G(f,g)
 &=-\int_0^L F(u)
 \frac{e^{-u/2}}{1-e^{-2u}}\,du.
\end{aligned}
\tag{20.2}
\]

\begin{theorem}[Complete source compensation]
\label{thm:source-compensation}
For every smooth-core Gate-SPG cross term,
\[
\boxed{
\begin{aligned}
 QW_L(f,g)
 &=
 -\int_0^L
 F(u)\frac{e^{-5u/2}}{1-e^{-2u}}\,du\\
 &\quad
 -\int_{(0,L]}F(u)e^{-u/2}\,
 d\bigl(\psi(e^u)-e^u\bigr).
\end{aligned}}
\tag{20.3}
\]
\end{theorem}

\begin{proof}
Since \(d(e^u)=e^u\,du\),
\[
 P+E
 =
 \int_0^L F(u)e^{-u/2}\,du
 -\int_{(0,L]}F(u)e^{-u/2}\,d(\Psi(u)-e^u).
\tag{20.4}
\]
The first term in (20.4) combines with Gamma through
\[
 e^{-u/2}-\frac{e^{-u/2}}{1-e^{-2u}}
 =-\frac{e^{-5u/2}}{1-e^{-2u}}.
\]
This proves (20.3).  The kernel is \(1/(2u)-3/4+O(u)\) near zero, and
\(F(0)=0\) makes the integral finite.
\end{proof}

The continuous PNT main measure has therefore cancelled the growing polar
branch exactly, while the remaining polar branch has combined with Gamma.
The sole unknown arithmetic distribution in (20.3) is
\(d(\psi(e^u)-e^u)\).  The explicit locally singular integral also has no
sign for a general cross correlation and must be controlled jointly.
Extension of the two separated integrals from the smooth core to arbitrary
Gate-SPG form-domain pairs requires a separate continuity argument.

\subsection*{The full quadratic completion}

For a zero-extended \(f\) put
\[
 H_f(u)=\langle f,\tau_uf\rangle,
 \qquad F_f(u)=H_f(u)+H_f(-u),
\]
and define
\[
\begin{aligned}
 \mathcal E_*(f)
 &:=
 \int_0^\infty\|f-\tau_uf\|_2^2
 \frac{e^{-5u/2}}{1-e^{-2u}}\,du,\\
 \mathcal A_\Delta(f)
 &:=
 \int_{(0,L]}F_f(u)e^{-u/2}\,
 d\bigl(\psi(e^u)-e^u\bigr),\\
 c_*&:=\gamma+\frac\pi2+3\log2+\log\pi-4>0.
\end{aligned}
\tag{20.5}
\]

\begin{theorem}[Exact quadratic prime--Gamma--pole compensation]
\label{thm:quadratic-compensation}
On the smooth form core,
\[
\boxed{
 QW_L(f,f)
 =
 \mathcal E_*(f)-c_*\|f\|_2^2-\mathcal A_\Delta(f).}
\tag{20.6}
\]
\end{theorem}

\begin{proof}
Let
\[
 d\nu_p(u)=e^{-u/2}d\psi(e^u),
 \quad d\nu_0(u)=e^{u/2}du,
 \quad d\nu_e(u)=e^{-u/2}du.
\]
Expanding the prime jump form in
Theorem~\ref{thm:centered-square} and subtracting its scalar part gives
\(-\int F_f\,d\nu_p+2\theta'(0)\|f\|^2\).  The growing polar branch cancels
\(\nu_0\), leaving
\[
 \int_0^LF_f\,d\nu_e-\mathcal A_\Delta(f).
\]
Since \(\nu_e(0,\infty)=2\),
\[
 \int_0^LF_f\,d\nu_e
 =4\|f\|_2^2-\int_0^\infty
 \|f-\tau_uf\|_2^2\,d\nu_e(u).
\]
Finally,
\[
 \frac{e^{-u/2}}{1-e^{-2u}}-e^{-u/2}
 =\frac{e^{-5u/2}}{1-e^{-2u}},
\]
and \(4+2\theta'(0)=-c_*\).  Combining the identities proves (20.6).
\end{proof}

Thus qualitative positivity of the completed form is exactly the
literal-prime inequality
\[
\boxed{
 \mathcal A_\Delta(f)
 \le\mathcal E_*(f)-c_*\|f\|_2^2
 \quad\text{for every test }f.}
\tag{20.7}
\]
Equation (20.7) is not proved here.  It is the one-sided signed theorem
left after the continuous PNT main term, Gamma and both polar branches have
been cancelled algebraically.

This reduction is already nontrivial on a normalized Fejer packet
\[
 f_{L,U}(x)=L^{-1/2}e^{iUx}\mathbf1_{I_L}(x).
\]
This is a diagonal calculation, with \(F(0)=2\), and is therefore not an
application of the cross identity (20.3).
Define
\[
\begin{aligned}
 S_L(U)&=\sum_{n\le e^L}\frac{\Lambda(n)}{\sqrt n}
 \left(1-\frac{\log n}{L}\right)\cos(U\log n),\\
 M_\pm(L,U)&=\int_0^Le^{\pm u/2}
 \left(1-\frac uL\right)\cos(Uu)\,du,\\
 \Delta_L(U)&=S_L(U)-M_+(L,U).
\end{aligned}
\tag{20.8}
\]
Put
\[
 A_+(L,U)=
 \frac{2\sinh((1/2+iU)L/2)}
 {\sqrt L\,(1/2+iU)}.
\]
Direct integration gives
\[
 P_L(U)
 =
 2\operatorname{Re}A_+(L,U)^2
 =2(M_++M_-),
\tag{20.9}
\]
or, equivalently,
\[
\begin{aligned}
 P_L(U)
 =\frac{4}{L(U^2+1/4)^2}\Big[
 &(1/4-U^2)(\cosh(L/2)\cos(UL)-1)\\
 &+U\sinh(L/2)\sin(UL)
 \Big].
\end{aligned}
\tag{20.9a}
\]
This polar value is not sign-definite.  Define the full archimedean packet
value by
\[
 A_{\infty,L}(U)
 =
 \int_{\mathbb R}2\theta'(t)
 |\widehat f_{L,U}(t)|^2\,\frac{dt}{2\pi}.
\tag{20.9b}
\]
Then
\[
\boxed{
 Q_L(f_{L,U},f_{L,U})
 =
 A_{\infty,L}(U)+2M_-(L,U)-2\Delta_L(U).}
\tag{20.10}
\]
The scalar signed target is therefore
\[
 \Delta_L(U)
 \le M_-(L,U)+\tfrac12A_{\infty,L}(U).
\tag{20.11}
\]
It is not proved here.

There is, however, an unconditional far-center result.  Put
\[
 c_\Gamma=\frac1{4\pi},
 \quad
 C_\Gamma=
 \sup_t\{c_\Gamma\log(2+|t|)-2\theta'(t)\},
 \quad
 C_\psi=\sup_{x\ge1}\frac{\psi(x)}x,
\]
and
\[
 m_0=\frac2\pi\int_0^1\left(\frac{\sin y}{y}\right)^2dy
 \ge\frac{2\sin^2(1)}{\pi}.
\]
The Fejer kernel of \(f_{L,U}\) has mass \(m_0\) on
\(|t-U|\le2/L\), while Stieltjes integration by parts gives
\[
 |S_L(U)|\le S_L(0)\le\frac{4C_\psi e^{L/2}}L.
\]
The exact polar formula (20.9a) also gives
\[
\begin{aligned}
 P_L(U)
 &\ge-\frac4L\left[
 \frac{\cosh(L/2)+1}{U^2+1/4}
 +\frac{|U|\sinh(L/2)}{(U^2+1/4)^2}
 \right]\\
 &\ge-\frac{5e^{L/2}}{LU^2}
 \qquad(L\ge2,\ |U|\ge4).
\end{aligned}
\]
It follows that, for \(L\ge2\) and \(|U|\ge4\),
\[
\boxed{
 Q_L(f_{L,U},f_{L,U})
 \ge
 m_0c_\Gamma\log(2+|U|/2)-C_\Gamma
 -\frac{8C_\psi e^{L/2}}L
 -\frac{5e^{L/2}}{LU^2}.}
\tag{20.12}
\]
Thus these packets are nonnegative once
\[
 \log(2+|U|/2)
 \ge
 \frac{C_\Gamma+(8C_\psi+5/16)\lambda/L}{m_0c_\Gamma}.
\tag{20.13}
\]
This improves the crude possible-center range from
\(\exp(C\lambda)\) to \(\exp(C\lambda/L)\) for this packet family.
It still leaves every fixed center as \(L\to\infty\).

Finally, exact polar tapering does not create second-moment center decay.
For
\[
 \mathcal P_L^\triangle(t)
 =
 \sum_{n\le e^L}\frac{\Lambda(n)}{\sqrt n}
 \left(1-\frac{\log n}{L}\right)n^{-it},
\]
The classical zero-free-region PNT and the two-sided
Montgomery--Vaughan theorem give, uniformly in \(U\),
\[
 \frac1{2Y}\int_{U-Y}^{U+Y}
 |\mathcal P_L^\triangle(t)|^2dt
 =
 \frac{L^2}{12}+O(1),
 \qquad Y=\frac{e^L}{L}.
\tag{20.14}
\]
The remaining theorem must therefore be genuinely one-sided and signed,
not another RMS estimate.

\section{The Picone gate and the positive-jump obstruction}
\label{sec:picone}

The positive square in Section~\ref{sec:centered-square} suggests a
nonlocal ground-state transform.  It exists exactly, but does not supply
the missing constant.

Write
\[
 \nu_N
 =
 \sum_{n\le N}\frac{\Lambda(n)}{\sqrt n}\delta_{\log n}
 +\frac{e^{-u/2}}{1-e^{-2u}}\,du.
\tag{21.1}
\]
Let \(v>0\) on \(I_L\), extend \(v\) by zero, and write \(r=f/v\) on
\(I_L\).

\begin{theorem}[Zero-extension Picone identity]
\label{thm:picone}
On a common smooth form core,
\[
\boxed{
\begin{aligned}
 \mathcal D_N(f,f)
 &=
 \int_{(0,\infty)}\!\int_{\mathbb R}
 v(x)v(x-u)|r(x)-r(x-u)|^2\,dx\,\nu_N(du)\\
 &\quad+
 \int_{I_L}\frac{(\mathcal D_Nv)(x)}{v(x)}
 |f(x)|^2\,dx.
\end{aligned}}
\tag{21.2}
\]
\end{theorem}

\begin{proof}
For each shift use
\[
 |v_xr_x-v_yr_y|^2
 =
 v_xv_y|r_x-r_y|^2
 +(v_x-v_y)(v_x|r_x|^2-v_y|r_y|^2).
\]
Set \(y=x-u\), integrate, and change variables in the second half of the
last term.  It becomes
\[
 \int_{I_L}
 \frac{2v(x)-v(x-u)-v(x+u)}{v(x)}|f(x)|^2\,dx.
\]
Integration against \(\nu_N\) proves (21.2), including the zero-extension
boundary killing.
\end{proof}

If \(\varphi_1>0\) is the normalized ground state of \(\mathcal D_N\),
with eigenvalue \(\mu_1\), then
\[
\begin{aligned}
 \mathcal D_N(f,f)-\mu_1\|f\|_2^2
 =
 \int\!\!\int
 \varphi_1(x)\varphi_1(x-u)
 \left|
 \frac{f(x)}{\varphi_1(x)}
 -\frac{f(x-u)}{\varphi_1(x-u)}
 \right|^2dx\,\nu_N(du).
\end{aligned}
\tag{21.3}
\]
Consequently \(\lambda_2(\mathcal D_N)\ge\kappa_N\) is exactly the
weighted Poincare inequality
\[
\int\!\!\int\varphi_1(x)\varphi_1(x-u)|r(x)-r(x-u)|^2
 dx\,\nu_N(du)
\ge(\kappa_N-\mu_1)\int\varphi_1^2|r|^2,
\tag{21.4}
\]
for \(\int\varphi_1^2r=0\).  The transform is lossless: its required
constant is the missing second-level bound itself.

The resolvent sign is equally exact.  Put
\(B_N=\mathcal D_N-\kappa_NI\), assume the strict one-negative-level
inertia from (19.14), and set
\(r_N=\langle h_L,B_N^{-1}h_L\rangle<0\).  Then
\[
 \inf_{\langle h_L,f\rangle=1}B_N(f,f)=\frac1{r_N},
\tag{21.5}
\]
so
\[
 1+2r_N\le0
 \quad\Longleftrightarrow\quad
 \inf_{\langle h_L,f\rangle=1}B_N(f,f)\ge-2.
\tag{21.6}
\]
Together with \(B_N\ge0\) on \(h_L^\perp\), this is precisely the original
completed even Weil inequality.

The following countermodel shows why generic jump positivity cannot prove
the inertia premise.

\begin{proposition}[Positive one-atom obstruction]
\label{prop:one-atom}
Fix \(a>0\) and \(M\ge6\).  There is a positive jump form with the same
Gamma channel, zero-extension boundary killing and polar rank-one
perturbation whose scalar-shifted compression has at least
\(\lfloor M/2\rfloor\) negative eigenvalues.
\end{proposition}

\begin{proof}
Choose normalized smooth bumps
\(\phi_j=\tau_{(j-1)a}\phi_1\), \(1\le j\le M\), with disjoint supports
of width below \(a/3\), inside a sufficiently long interval.  Add the
positive atom \(w\delta_a\) and subtract its scalar contribution \(2wI\),
as in (19.3).  On their span the atomic part is
\[
 -wA_M,
\]
where \(A_M\) is the path adjacency matrix, with eigenvalues
\[
 -2w\cos\frac{k\pi}{M+1},
 \qquad1\le k\le M.
\]
The fixed Gamma, scalar and polar compressions have norm independent of
\(w\).  Weyl's inequality therefore leaves at least
\(\lfloor M/2\rfloor\) negative levels for sufficiently large \(w\).
\end{proof}

Thus positivity of the jump measure, Picone algebra, Hardy completion and
rank-one perturbation do not select the fundamental branch.  The actual
locations and magnitudes \(\Lambda(n)/\sqrt n\) must enter through the
literal compensated inequality (20.7), quantitatively strengthened to the
moving-vector residual and \(d_8\) scale of Gate SPG.

\section{Smooth full-kernel operator quasimodes}
\label{sec:full-kernel-quasimode}

The moving co-Poisson construction retains the finite-cutoff geometry of
CCM.  For the direct positivity problem there is a second route: truncate
Riemann's full kernel only after using its exact radical property.

Let \(K\) be the even positive Schwartz function normalized by
\[
 \widehat K(z)=\Xi(z).
\tag{22.1}
\]
Write \(a=L/2\), \(\lambda=e^a\), and choose smooth even cutoffs
\(\chi_a\) equal to one on \([-a+1,a-1]\) and supported in
\((-a,a)\), with all fixed derivative seminorms bounded uniformly in
\(a\).  Put
\[
 v_L=\chi_aK,\qquad q_L^K=\frac{v_L}{\|v_L\|_2}.
\tag{22.2}
\]

The theta expansion and modular symmetry imply, for every fixed \(j\),
\[
 |K^{(j)}(x)|
 \le C_j e^{M_j|x|}e^{-c_je^{2|x|}}.
\tag{22.3}
\]
Consequently, if \(e_L=v_L-K\), repeated integration by parts gives
\[
 \sup_{|\eta|\le1/2}|\widehat e_L(t+i\eta)|
 \le C_r\lambda^{M_r}e^{-c_r\lambda^2}(1+|t|)^{-r}.
\tag{22.4}
\]

\begin{theorem}[Full-kernel operator quasimode]
\label{thm:full-kernel-quasimode}
For suitable positive constants \(C,M,c\),
\[
 \|A_L^+v_L\|_2
 \le C\sqrt{L\lambda}\,\lambda^Me^{-c\lambda^2}.
\tag{22.5}
\]
In particular,
\[
 \|A_L^+q_L^K\|_2\longrightarrow0.
\tag{22.6}
\]
\end{theorem}

\begin{proof}
At every nontrivial zero \(z_\rho\) in the additive spectral coordinate,
\(\widehat K(z_\rho)=0\).  Thus the polarized Weil form satisfies
\[
 QW(K,g)=0,
\tag{22.7}
\]
and \(QW(v_L,g)=QW(e_L,g)\).  For \(g\in L^2([-a,a])\),
\[
 |\widehat g(z_\rho)|
 \le\sqrt{L\lambda}\,\|g\|_2,
\tag{22.8}
\]
because \(|\operatorname{Im}z_\rho|<1/2\).  The local
Riemann--von Mangoldt estimate gives
\(\sum_\rho(1+|\gamma|)^{-r}<\infty\) for \(r>1\).  Combining
(22.4), the polarized zero sum and (22.8) proves
\[
 |QW(v_L,g)|
 \le C\sqrt{L\lambda}\,\lambda^Me^{-c\lambda^2}\|g\|_2.
\]
The Riesz representative of this functional is \(A_L^+v_L\).  Finally
\(v_L\to K\) in \(L^2\), so its norm stays bounded away from zero.
\end{proof}

Theorem~\ref{thm:full-kernel-quasimode} closes the residual half of the
finite-cutoff branch criterion.  It does not supply its complementary
spectral floor.  Ordinary cyclicity of translates and derivatives of
\(K\) cannot supply that floor at the required scale.  Indeed, if an
off-line zero \(\rho=\beta+i\gamma\), \(\beta>1/2\), exists and a
fixed compactly supported \(h\) detects it, then every truncated radical
approximant \(w_L\) with
\(\widehat w_L(z_\rho)=o(1)\) satisfies
\[
 \|h-w_L\|_2^2\gg_{h,\rho}
 \lambda^{-2(\beta-1/2)},\qquad
 \kappa_{e^L}\|h-w_L\|_2^2
 \gg_{h,\rho}\lambda^{2(1-\beta)}\longrightarrow\infty,
\tag{22.9}
\]
where \(\kappa_{e^L}=4\lambda+o(\lambda)\).  Thus qualitative
\(L^2\)-density has exactly the wrong conditioning for the semilocal
inertia problem.

\section{A cutoff-free normalization of the classical Weil criterion}
\label{sec:doob-variance}

Weil positivity is a classical criterion equivalent to RH
\cite{Weil1952}.  Many later works give other analytic, Hilbert-space and
dynamical equivalences; in particular Sun derives an explicit integral
inequality from \(|\xi(\sigma-it)|^2\) and an equivalent infinite
orthogonalization-time statement on Wiener space \cite{Sun}.  The purpose
of this section is not to add another logical equivalent.  It is to express
the same classical Weil form in the exact ground coordinate needed to
compare it with the finite CCM--Suzuki operators and with the literal
prime-power decomposition of the preceding sections.

The full kernel removes the moving branch from the algebra before any
estimate is made.  Put
\[
 h(x)=\cosh(x/2),\qquad
 c_K=\int_\mathbb R h(x)K(x)\,dx
     =\widehat K(i/2)=\xi(0)=\frac12
\tag{23.1}
\]
and define the probability measure
\[
 d\mu_K(x)=\frac{h(x)K(x)}{c_K}\,dx
          =2\cosh(x/2)K(x)\,dx.
\tag{23.2}
\]
The complete positive jump measure is
\[
 \nu_\zeta(du)
 =\sum_{n\ge2}\frac{\Lambda(n)}{\sqrt n}
       \delta_{\log n}(du)
  +\frac{e^{-u/2}}{1-e^{-2u}}\,du,
 \qquad u>0.
\tag{23.3}
\]
For even \(r\in C_c^\infty(\mathbb R)\), set
\[
 \mathscr E_K(r)
 =\int_{(0,\infty)}\!\int_\mathbb R
 K(x)K(x-u)|r(x)-r(x-u)|^2\,dx\,\nu_\zeta(du).
\tag{23.4}
\]

\begin{theorem}[Full-kernel Doob--variance identity]
\label{thm:doob-variance}
For every even complex \(r\in C_c^\infty(\mathbb R)\),
\[
 \boxed{
 QW(Kr,Kr)=\mathscr E_K(r)
 -\frac12\operatorname{Var}_{\mu_K}(r).}
\tag{23.5}
\]
\end{theorem}

\begin{proof}
Use a common smooth spatial regulator and the same prime cutoff in the two
physical-side Weil forms \(QW(Kr,Kr)\) and
\(QW(K,K|r|^2)\).  At every retained jump, with
\(y=x-u\), the pointwise identity
\[
 |K_xr_x-K_yr_y|^2
 -(K_x-K_y)(K_x|r_x|^2-K_y|r_y|^2)
 =K_xK_y|r_x-r_y|^2
\tag{23.6}
\]
gives (23.4).  The scalar centering terms cancel.  In the even sector the
two polar evaluations contribute
\[
 2\left|\int hKr\right|^2
 -2\left(\int hK\right)
    \left(\int hK|r|^2\right)
 =-2c_K^2\operatorname{Var}_{\mu_K}(r).
\tag{23.7}
\]
Since \(2c_K^2=1/2\), this is the second term of (23.5).

The zero-side radical identity \(QW(K,g)=0\), applied to
\(g=K|r|^2\), removes the second regulated Weil form in the limit.
The Gamma contribution is integrable at zero because the squared
difference in (23.4) is \(O(u^2)\), and the prime sum converges after the
spatial regulator is removed because one translated \(K\)-factor has
double-exponential decay.  Monotone convergence applies to the nonnegative
prime energy.  Removing the two regulators therefore proves (23.5).
\end{proof}

The jump form (23.4) is closable on \(L^2(\mu_K)\).  Indeed, define the
positive symmetric \(\sigma\)-finite measure \(J_K\) by
\[
 \int F(x,y)\,J_K(dx,dy)
 =\frac12\int_{(0,\infty)}\!\int_\mathbb R
 K(x)K(x-u)
 \bigl(F(x,x-u)+F(x-u,x)\bigr)\,dx\,\nu_\zeta(du).
\tag{23.7a}
\]
Then (23.4) is the standard jump form
\(\int|r(x)-r(y)|^2J_K(dx,dy)\), initially on its smooth core.
Let
\(\mathcal L_K\) denote the associated nonnegative self-adjoint operator,
whose weak action on the core is
\[
 (\mathcal L_Kr)(x)
 =\frac{c_K}{h(x)}
 \int_{(0,\infty)}
 \begin{aligned}[t]
 \bigl[&K(x-u)(r(x)-r(x-u))\\
       &+K(x+u)(r(x)-r(x+u))\bigr]\,\nu_\zeta(du).
 \end{aligned}
\tag{23.7b}
\]
Let
and put
\[
 \mathcal H_{0,+}
 =\left\{r\in L^2(\mu_K):
 r(-x)=r(x),\ \int r\,d\mu_K=0\right\}.
\tag{23.8}
\]

\begin{corollary}[Exact spectral floor]
\label{cor:exact-spectral-floor}
The following statements are equivalent:
\[
 \mathrm{RH};
 \qquad
 \mathscr E_K(r)\ge\frac12\operatorname{Var}_{\mu_K}(r)
 \quad(r\in C_{c,\mathrm{even}}^\infty);
 \qquad
 \inf\sigma(\mathcal L_K|_{\mathcal H_{0,+}})\ge\frac12.
\tag{23.9}
\]
\end{corollary}

\begin{proof}
Every compactly supported smooth even test \(f\) has the unique form
\(f=Kr\), since \(K>0\).  By the bilateral parity theorem, failure of
RH already produces a negative even Weil test.  Theorem
\ref{thm:doob-variance} therefore converts the even Weil criterion exactly
into the middle inequality of (23.9).  Closing the form gives the spectral
statement.
\end{proof}

Corollary~\ref{cor:exact-spectral-floor} is a reformulation of the classical
Weil criterion.  It is not, by itself, progress toward proving the missing
positivity, nor does it establish the strong-resolvent convergence proposed
in the CCM--Suzuki program.  Its role here is bookkeeping: it places the
remaining sign question in the same full-kernel coordinate as the exact
theta and von Mangoldt decompositions below.

The threshold is not an arbitrary estimate.  It already contains an
infinite radical family.

\begin{proposition}[Infinite threshold family]
\label{prop:threshold-family}
For \(m\ge1\), define
\[
 r_m=\frac{K^{(2m)}}K-4^{-m}.
\tag{23.10}
\]
Then \(r_m\in\mathcal H_{0,+}\) belongs to the extended form domain and
is a weak eigenfunction of \(\mathcal L_K\) with eigenvalue \(1/2\).
The family \((r_m)_{m\ge1}\) is linearly independent.  Consequently
\(1/2\) is an infinite-multiplicity, hence essential, spectral point.
\end{proposition}

\begin{proof}
Integration by parts and \(h^{(2m)}=4^{-m}h\) give
\[
 \int\frac{K^{(2m)}}K\,d\mu_K=4^{-m},
\]
so \(r_m\) has mean zero.  Moreover
\[
 \widehat{K^{(2m)}}(z)=(-1)^mz^{2m}\Xi(z),
\]
and both \(K^{(2m)}\) and \(K\) lie in the polarized Weil radical.
Polarizing (23.5) therefore yields
\[
 \mathscr E_K(r_m,s)
 =\frac12\langle r_m,s\rangle_{L^2(\mu_K)}
\tag{23.11}
\]
for every even test \(s\), and cutoff approximation extends the identity
to the form domain.  Finally, the theta asymptotics make
\(K^{(2m)}/K\) a function with a distinct highest growth degree in
\(e^{2|x|}\) for each \(m\); hence no nontrivial finite linear
combination can vanish.
\end{proof}

Proposition~\ref{prop:threshold-family} supplies an important qualification.
The currently proved equivalence is the absence of \emph{any} spectrum below
\(1/2\).  Calling every possible subthreshold point a bound state requires
one additional structural theorem,
\[
 \inf\sigma_{\mathrm{ess}}
 (\mathcal L_K|_{\mathcal H_{0,+}})\ge\frac12,
\tag{23.12}
\]
after which the remaining exclusion would indeed concern only discrete
bound states.

\section{Atomwise rigidity and the all-source sharp gate}
\label{sec:atomwise-rigidity}

For \(u>0\), write
\[
 \mathcal J_u(r)
 =\int_\mathbb R K(x)K(x-u)
 |r(x)-r(x-u)|^2\,dx.
\tag{24.1}
\]
Then
\[
 \mathscr E_K(r)
 =\mathscr E_\Gamma(r)
  +\sum_{m\ge2}\frac{\Lambda(m)}{\sqrt m}
       \mathcal J_{\log m}(r).
\tag{24.2}
\]
This is an increasing hierarchy in the set of retained prime-power atoms.
Positivity of the omitted tail, however, cannot be used to reduce the
problem to finitely many atoms.

\begin{theorem}[Every von Mangoldt atom is load-bearing]
\label{thm:atomwise-load}
Let \(S\) omit at least one prime power, and define
\[
 \mathscr E_S(r)
 =\mathscr E_\Gamma(r)
  +\sum_{m\in S}\frac{\Lambda(m)}{\sqrt m}
       \mathcal J_{\log m}(r).
\tag{24.3}
\]
Then the sharp inequality
\(\mathscr E_S(r)\ge\tfrac12\operatorname{Var}_{\mu_K}(r)\)
is false.  More precisely, there are even compactly supported smooth
\(r_a\) for which the reverse strict inequality holds for every
sufficiently large \(a\).
\end{theorem}

\begin{proof}
Take
\[
 f_a=\chi_aK'',\qquad
 r_a=\frac{f_a}{K}.
\tag{24.4}
\]
The transform \(\widehat{K''}(z)=-z^2\Xi(z)\) and the exterior estimate
of Theorem~\ref{thm:full-kernel-quasimode} give
\(QW(f_a,f_a)\to0\).  For every fixed \(u>0\),
\[
 \mathcal J_u(r_a)\longrightarrow
 \mathcal J_u(K''/K)>0.
\tag{24.5}
\]
Strict positivity follows because equality would make \(K''/K\)
\(u\)-periodic, whereas its theta asymptotic is unbounded.  Subtracting the
omitted atoms from (23.5), and retaining one omitted \(m_0\), gives
\[
 \limsup_{a\to\infty}
 \left(
 \mathscr E_S(r_a)-\frac12\operatorname{Var}_{\mu_K}(r_a)
 \right)
 \le
 -\frac{\Lambda(m_0)}{\sqrt{m_0}}
 \mathcal J_{\log m_0}(K''/K)<0.
\tag{24.6}
\]
\end{proof}

The theta series contains an exact dilation which exposes the missing
coupling.  Suppressing the common positive normalization, put
\[
 k_j(x)=\pi j^2e^{5x/2}
 \left(2\pi j^2e^{2x}-3\right)e^{-\pi j^2e^{2x}},
 \qquad K(x)=\sum_{j\ge1}k_j(x)
 \quad(x\ge0).
\tag{24.7}
\]
For every integer \(n\ge2\),
\[
 \boxed{k_{nj}(x-\log n)=n^{-1/2}k_j(x).}
\tag{24.8}
\]
Consequently, when \(x\ge\log n\),
\[
 K(x-\log n)
 =n^{-1/2}K(x)
 +\sum_{\substack{j\ge1\\n\nmid j}}
   k_j(x-\log n).
\tag{24.9}
\]
Dropping the second term and the central crossing region
\(0<x<\log n\) produces a valid lower bound, but Theorem
\ref{thm:atomwise-load} shows that it is strictly too small at the sharp
constant.  Thus the all-prime theorem must preserve simultaneously:
\begin{enumerate}
\item every atom \(\Lambda(p^k)\);
\item the theta indices divisible by the corresponding dilation;
\item the nondivisible remainder in (24.9);
\item the central crossing region;
\item the Gamma jump continuum and the polar variance.
\end{enumerate}

The arithmetic identity
\[
 \sum_{d\mid j}\Lambda(d)=\log j
\tag{24.10}
\]
shows the natural point at which the divisible theta channel can couple to
the logarithmic derivative implicit in the Gamma term.  It does not control
the nondivisible and crossing channels separately; those channels cannot be
discarded, by (24.6).

There is, however, an exact reorganization which retains all of those
channels.  Write
\(\mathscr E_p=\mathscr E_K-\mathscr E_\Gamma\).

\begin{proposition}[Exact half-line theta--divisor lift]
\label{prop:theta-divisor-lift}
For every even \(r\) in the test core,
\[
 \mathscr E_p(r)
 =\mathscr E_p^{\rm same}(r)+\mathscr E_p^{\rm cross}(r),
\tag{24.11}
\]
where
\[
 \begin{aligned}
 \mathscr E_p^{\rm same}(r)
 &=2\int_0^\infty K(x)
   \sum_{\ell\ge2}k_\ell(x)
   \sum_{\substack{d\mid\ell\\d\ge2}}
   \Lambda(d)|r(x)-r(x+\log d)|^2\,dx,\\
 \mathscr E_p^{\rm cross}(r)
 &=\sum_{d\ge2}\frac{\Lambda(d)}{\sqrt d}
   \int_0^{\log d}K(t)K(\log d-t)
   |r(t)-r(\log d-t)|^2\,dt.
 \end{aligned}
\tag{24.12}
\]
The Gamma energy has the parallel decomposition
\[
 \begin{aligned}
 \mathscr E_\Gamma^{\rm same}(r)
 &=2\int_0^\infty\!\int_0^\infty
 K(x)K(x+u)|r(x)-r(x+u)|^2\,dx\,\nu_\Gamma(du),\\
 \mathscr E_\Gamma^{\rm cross}(r)
 &=\int_0^\infty\!\int_0^u
 K(t)K(u-t)|r(t)-r(u-t)|^2\,dt\,\nu_\Gamma(du),
 \end{aligned}
\tag{24.13}
\]
with \(\nu_\Gamma(du)=e^{-u/2}(1-e^{-2u})^{-1}du\).
Moreover, if
\[
 d\mu_K^+(x)=2\mathbf1_{[0,\infty)}(x)d\mu_K(x),
\tag{24.14}
\]
then \(\mu_K^+\) is a probability measure and
\[
 \operatorname{Var}_{\mu_K}(r)
 =\operatorname{Var}_{\mu_K^+}(r).
\tag{24.15}
\]
\end{proposition}

\begin{proof}
Orient the prime displacement as
\[
 \int_\mathbb R K(x)K(x+\log d)
 |r(x)-r(x+\log d)|^2\,dx.
\]
The two same-sign rays \(x\ge0\) and
\(x\le-\log d\) agree by evenness and give twice the positive
half-line integral.  The interval \((-\log d,0)\), after reflection,
gives the second line of (24.12).  From (24.8),
\[
 k_m(x+\log d)=\sqrt d\,k_{dm}(x).
\tag{24.16}
\]
Therefore
\[
 \frac{\Lambda(d)}{\sqrt d}K(x)K(x+\log d)
 =\Lambda(d)K(x)\sum_{m\ge1}k_{dm}(x).
\]
Summing in \(d\) and grouping \(\ell=dm\) proves the first line of
(24.12).  The same ray/crossing split gives (24.13).  Finally, \(K\),
\(h\) and \(r\) are even, so conditioning \(\mu_K\) on the positive
half-line changes neither moment of \(r\), proving (24.15).
\end{proof}

Proposition~\ref{prop:theta-divisor-lift} is not another RH-equivalent
criterion.  It is an exact arithmetic coordinate for the already-known
criterion: a half-line chain with theta--divisor outward edges, central
reflection edges and Gamma continuum edges.  The outer factor \(K(x)\)
retains all theta indices, including those which are nondivisible for a
given displacement; the crossing channel is kept rather than estimated
away.

\begin{problem}[All-source sharp gate]
\label{prob:all-source-gate}
Prove, with all terms in (24.2) retained,
\[
 \boxed{
 \mathscr E_K(r)\ge
 \frac12\operatorname{Var}_{\mu_K}(r)
 \qquad
 (r\in C_{c,\mathrm{even}}^\infty(\mathbb R)).}
\tag{24.17}
\]
Equivalently, prove
\[
 \sigma(\mathcal L_K|_{\mathcal H_{0,+}})
 \subset[1/2,\infty).
\tag{24.18}
\]
\end{problem}

A constructive closure of (24.17) may take either of two equivalent
operator forms.  The strongest is an arithmetic factorization
\[
 \mathscr E_K(r)-\frac12\operatorname{Var}_{\mu_K}(r)
 =\|\mathcal Br\|_{\mathcal Y}^2
\tag{24.19}
\]
on the quotient by the threshold radical, where \(\mathcal B\) must be
built from the complete divisor--theta coupling rather than from generic
jump positivity.  A second route is to prove (23.12), choose a reference
operator with threshold \(1/2\), and exclude its subthreshold
Birman--Schwinger spectrum by an all-source norm bound.  In either route,
the equality modes of Proposition~\ref{prop:threshold-family} determine the
sharp normalization and forbid any strict global improvement of the
constant.

More explicitly, put
\[
 \mathcal R_K=
 \overline{\operatorname{span}}
 \left\{K^{(2m)}/K-4^{-m}:m\ge1\right\}
 \subset\mathcal H_{0,+}.
\tag{24.20}
\]
Proposition~\ref{prop:threshold-family} gives
\(\mathcal L_K|_{\mathcal R_K}=\tfrac12I\).  The unresolved
coercive statement is therefore the same sharp floor on
\(\mathcal R_K^\perp\), with the three exact channels
(24.12)--(24.13) kept coupled.

\section{Conclusion}

Starting with the finite self-adjoint quotients of CCM and the continuous
localized perspective of Suzuki, this paper constructs a single chain
through the two limiting parameters.  At fixed physical cutoff, the
quotients possess a canonical metric-compatible transport; its defect is a
rank-two source expression plus its Loewner metric correction.  This yields
the resolvent-shell identities, cofinal fixed-cutoff summability and the
precise adjacent-cutoff obstruction.

In the outer limit, logarithmic curvature is shown to be an affine-rigid
divisor observable.  The moving co-Poisson construction keeps primes,
prime powers, Gamma and pole coupled; the centered jump identity restores
the Fourier center lost by translated mean-square estimates; and exact PNT
compensation isolates the signed discrepancy channel.  These results
identify what the finite operator convergence program must still prove:
the global branch selector together with the moving/model identification.

The full Riemann kernel provides a second, cutoff-free endpoint of the same
development.  Its smooth truncations are operator quasimodes, and the exact
radical identity converts the completed Weil form into
\[
 QW(Kr,Kr)
 =\mathscr E_K(r)-\frac12\operatorname{Var}_{\mu_K}(r).
\]
The resulting RH equivalence is a ground-state normalization of the
classical Weil criterion, not a new criterion and not progress by itself
toward proving its sign.  Its role in this paper is to expose the literal
prime-power and theta channels of the unresolved global branch.
Within that normalization, the polar term gives the exact constant
\(1/2\), and the complete
prime-power--Gamma measure remains visible in \(\mathscr E_K\).
Corollary~\ref{cor:exact-spectral-floor} therefore identifies direct Weil
positivity, hence RH, with the spectral floor
\[
 \inf\sigma(\mathcal L_K|_{\mathcal H_{0,+}})\ge\frac12.
\]

The threshold is sharp and highly degenerate:
Proposition~\ref{prop:threshold-family} places an infinite family of
Riemann-radical modes exactly at \(1/2\).  Theorem
\ref{thm:atomwise-load} also proves that every literal von Mangoldt atom is
load-bearing.  The exact theta dilation cannot be weakened to its divisible
part, because the nondivisible theta indices and the central crossing region
carry a strictly positive part of the threshold energy.  Thus finite-prime
certificates, positive-tail completion, generic Picone identities and
termwise dilation estimates cannot establish the sharp floor.

There are consequently two related but distinct unresolved statements.
For convergence of the particular CCM--Suzuki spectral approximants, the
remaining conditions are the signed branch selector and moving/model
identification.  For a direct proof of RH in the full-kernel coordinate, the
single remaining theorem is Problem~\ref{prob:all-source-gate}: preserve
all prime atoms, divisible and nondivisible theta indices, the central
crossing, Gamma and the polar variance, and prove that no spectrum lies
below \(1/2\).  Proving first that the essential spectrum has floor
\(1/2\) would reduce the final step to excluding discrete subthreshold
states.  Neither the sharp inequality nor the full operator convergence is
claimed here.

\begin{thebibliography}{99}

\bibitem{Weil1952}
A.~Weil,
\emph{Sur les ``formules explicites'' de la théorie des nombres premiers},
Comm. Sém. Math. Univ. Lund (1952), 252--265.

\bibitem{Sun}
W.~Sun,
\emph{An equivalent inequality for the Riemann hypothesis},
arXiv:2312.12760v5 (2024).

\bibitem{CCM}
A.~Connes, C.~Consani, and H.~Moscovici,
\emph{Zeta spectral triples},
arXiv:2511.22755 (2025).

\bibitem{Suzuki}
M.~Suzuki,
\emph{Weil's quadratic form via the screw function},
arXiv:2606.09096 (2026).

\bibitem{Connes2026}
A.~Connes,
\emph{The Riemann Hypothesis: Past, Present and a Letter Through Time},
arXiv:2602.04022 (2026).

\bibitem{MV}
H.~L.~Montgomery and R.~C.~Vaughan,
\emph{Hilbert's inequality},
Journal of the London Mathematical Society (2) \textbf{8} (1974), 73--82.

\bibitem{Bhatia}
R.~Bhatia,
\emph{Matrix Analysis},
Graduate Texts in Mathematics 169, Springer, 1997.

\bibitem{Kato}
T.~Kato,
\emph{Perturbation Theory for Linear Operators},
Classics in Mathematics, Springer, 1995.

\end{thebibliography}

\end{document}
